\chapter{Digital Filters and Multirate Processing}
\label{ch:dsp}

\begin{chapterintro}
A software-defined radio is, to a first approximation, a collection of digital filters
running at different sample rates. Channel selection, pulse shaping, matched filtering,
interpolation for timing recovery, decimation after the ADC and interpolation before the DAC
are all filtering problems, and the efficiency with which they are solved determines whether
a design fits in a phone, an FPGA or a satellite. This chapter is the working engineer's
guide to digital filters. It starts with the $z$-transform as a tool for intuition rather
than for contour integrals: poles, zeros and the geometry that turns a pole-zero plot into a
frequency response. It then covers how to specify a filter and estimate its cost before
designing it; the FIR design methods (windows, frequency sampling, least squares and the
Parks--McClellan equiripple algorithm) and the special filters every radio uses (half-band,
Nyquist, Hilbert); the classical IIR families and the bilinear transform; and the
implementation questions that decide whether a design works in silicon: structures,
coefficient quantization, limit cycles, overflow, scaling and fixed-point word lengths. The
second half is about changing the sample rate efficiently: decimation and interpolation, the
noble identities, polyphase decomposition, multistage design, half-band chains,
Hogenauer's multiplier-free CIC filter, the polyphase DFT filter bank that channelizes a
satellite payload, and arbitrary resampling. It closes with the oscillators and rotation
engines---the NCO and CORDIC---and assembles everything into the digital down- and
up-converters inside every modern radio, from a GC4016 to the AD9361 in a USRP B200.
\end{chapterintro}

%======================================================================
\section{Discrete-time systems and the \texorpdfstring{$z$}{z}-transform}
\label{sec:ch06:ztrans}
%======================================================================

Chapter~\ref{ch:signals} introduced linear time-invariant (LTI) systems, the DTFT and the DFT.
This section adds the one tool that chapter used only in passing and that a filter designer
uses every day: the $z$-transform, read not as a branch of complex analysis but as a picture.

\subsection{Difference equations and transfer functions}

A discrete-time LTI system is characterised by its impulse response $h[n]$ and computes the
convolution $y[n]=\sum_k h[k]\,x[n-k]$. Its \term{z-transform}
\begin{equation}
H(z)=\sum_{n=-\infty}^{\infty}h[n]\,z^{-n}
\label{eq:ch06:ztrans}
\end{equation}
turns convolution into multiplication, $Y(z)=H(z)X(z)$, and its frequency response is
$H(e^{j\omega})$, the $z$-transform evaluated on the unit circle, with $\omega=2\pi f/f_s$ in
radians per sample. The only property of $z^{-1}$ you need is that it is a one-sample delay: if
$x[n]$ has transform $X(z)$, then $x[n-1]$ has transform $z^{-1}X(z)$. Every block diagram in this
chapter is drawn with $z^{-1}$ boxes for that reason.

Almost every practical filter is described by a linear constant-coefficient difference equation,
\begin{equation}
\begin{aligned}
y[n]&=\sum_{k=0}^{M}b_k\,x[n-k]-\sum_{k=1}^{N}a_k\,y[n-k]\\
\Longleftrightarrow\quad
H(z)&=\frac{\sum_{k=0}^{M}b_kz^{-k}}{1+\sum_{k=1}^{N}a_kz^{-k}}
=b_0\,z^{N-M}\frac{\prod_{k=1}^{M}(z-z_k)}{\prod_{k=1}^{N}(z-p_k)}.
\end{aligned}
\label{eq:ch06:diffeq}
\end{equation}
The roots $z_k$ of the numerator are the \term{zeros} and the roots $p_k$ of the denominator the
\term{poles}. With real coefficients they are real or occur in complex-conjugate pairs. If
$a_k=0$ for all $k$, the filter has only zeros (apart from poles at the origin, which are pure
delays) and a finite impulse response: it is an \term{FIR} filter. Otherwise the feedback makes
the impulse response infinitely long: an \term{IIR} filter. Choosing between them is the first
decision in any filter design, and most of this chapter is about making it well.

\subsection{Stability and the region of convergence}

The sum~\eqref{eq:ch06:ztrans} converges only for some values of $z$, its \emph{region of
convergence} (ROC). For a causal sequence ($h[n]=0$ for $n<0$) the ROC is the outside of a circle
that passes through the outermost pole. A system is bounded-input bounded-output (BIBO) stable if
and only if $\sum_n|h[n]|<\infty$, which is the same as requiring the unit circle to lie inside the
ROC. Combining the two statements gives the rule every engineer memorises: \emph{a causal system
is stable if and only if all its poles lie strictly inside the unit circle}. A pole at radius $r$
contributes a term $r^n$ to the impulse response, so $r=0.99$ decays by a factor $e$ in about 100
samples and $r=0.999$ in about 1000: the closer a pole is to the circle, the longer the filter
rings, the narrower the feature it produces in the frequency response, and---as
Section~\ref{sec:ch06:iirnum} shows---the more sensitive it is to arithmetic errors. FIR filters
have no poles outside the origin and are unconditionally stable, which is one of their great
practical virtues.

\subsection{Frequency response from pole-zero geometry}

Evaluate the factored form of~\eqref{eq:ch06:diffeq} at a point $z=e^{j\omega}$ on the unit circle.
Each factor $e^{j\omega}-z_k$ is a vector from the zero $z_k$ to that point, and each
$e^{j\omega}-p_k$ a vector from a pole. Therefore
\begin{equation}
\begin{aligned}
|H(e^{j\omega})|&=|b_0|\,\frac{\prod_k|e^{j\omega}-z_k|}{\prod_k|e^{j\omega}-p_k|},\\
\angle H(e^{j\omega})&=\angle b_0+\sum_k\angle(e^{j\omega}-z_k)-\sum_k\angle(e^{j\omega}-p_k)+(N-M)\omega.
\end{aligned}
\label{eq:ch06:geometry}
\end{equation}
As $\omega$ sweeps from $0$ to $\pi$, the point travels counter-clockwise around the upper half of
the circle. When it passes close to a zero, one vector in the numerator becomes short and the
response dips---to zero if the zero lies on the circle. When it passes close to a pole, a vector in
the denominator becomes short and the response peaks. That is all there is to it, and with a
little practice you can sketch the magnitude response of any low-order filter from its pole-zero
plot in a few seconds. Every experienced DSP engineer does.

\mdcfig{ch06_pz_geometry}{Reading a frequency response from the pole-zero plot. Left: a notch filter
with zeros on the unit circle at $\pm0.15$ cycles/sample and poles just inside at radius $r$. The
magnitude response at the point $e^{j\omega}$ is the product of the zero-vector lengths divided by
the product of the pole-vector lengths. Right: far from the notch the two vectors of each pair are
almost equal and the response is flat; near it, the zero vector shrinks to nothing. The closer the
poles are to the circle, the narrower the notch.}

Figure~\ref{fig:ch06_pz_geometry} shows the idea at work in a \term{notch filter}, the classic
two-pole, two-zero section
\begin{equation}
H(z)=\frac{1-2\cos\omega_0\,z^{-1}+z^{-2}}{1-2r\cos\omega_0\,z^{-1}+r^2z^{-2}}.
\label{eq:ch06:notch}
\end{equation}
The zeros sit on the circle at $e^{\pm j\omega_0}$, so the response is exactly zero at $\omega_0$.
The poles sit at the same angles but at radius $r<1$. Away from $\omega_0$ each pole vector is
almost the same length as its zero vector, so they cancel and the response is close to unity.
Only within a distance of about $1-r$ of the zero does the zero vector become much shorter than the
pole vector. The $-3$\,dB width of the notch is therefore about $2(1-r)$ radians per sample, or
$(1-r)/\pi$ cycles per sample, a result worth remembering because it applies to any isolated
pole: a pole at radius $r$ produces a resonance whose $-3$\,dB bandwidth is about $(1-r)f_s/\pi$\,Hz.

\begin{worked}[A notch for a narrowband interferer]
A receiver sampled at $f_s=48$\,kHz suffers from a spur at 1\,kHz, and the designer wants to remove
it with a notch whose $-3$\,dB width is 20\,Hz. The notch frequency is $\omega_0=2\pi\cdot1000/48000=
0.1309$\,rad/sample, so $2\cos\omega_0=1.9829$. The width condition gives
$1-r=\pi\cdot20/48000=1.31\times10^{-3}$, so $r=0.99869$, and~\eqref{eq:ch06:notch} becomes
\[
H(z)=\frac{1-1.9829z^{-1}+z^{-2}}{1-1.9803z^{-1}+0.99738z^{-2}}.
\]
Five multiplications per sample remove the spur and leave the rest of the band within a small
fraction of a decibel. The poles at radius 0.9987 also tell you the cost: the filter's impulse
response decays as $0.9987^n$, a time constant of about 760 samples or 16\,ms, so the notch
needs that long to settle when the interferer appears. Narrow in frequency means long in time,
whatever the structure. Note also that the denominator coefficients differ from the numerator's
in the fourth significant figure; Section~\ref{sec:ch06:iirnum} explains why such filters need
care in fixed point.
\end{worked}

The same reasoning designs a \emph{resonator} (poles near the circle, zeros at $z=\pm1$ to suppress
DC and $f_s/2$), the \emph{DC blocker} $H(z)=(1-z^{-1})/(1-\alpha z^{-1})$ that removes the
residual offset of every direct-conversion receiver (a zero at $z=1$ with a pole just inside it),
and the \emph{comb filter} $1-z^{-R}$, whose $R$ zeros are spread evenly around the circle. The
comb filter will reappear as half of the CIC decimator in Section~\ref{sec:ch06:cic}.

\subsection{Minimum phase, maximum phase and all-pass}

Chapter~\ref{ch:signals} showed that the magnitude response does not determine the phase. The
$z$-plane makes the reason visible. Reflect a zero $z_k$ to its conjugate reciprocal $1/z_k^*$:
by elementary geometry the distance from every point on the unit circle to $1/z_k^*$ is
$1/|z_k|$ times the distance to $z_k$, so the magnitude response changes only by a constant while
the phase changes. A filter can therefore have any of $2^M$ zero configurations with the same
magnitude shape. The one with all zeros inside the circle is \term{minimum phase}: it has the
smallest phase lag and the most compact energy distribution of the family, and it has a stable
causal inverse. The one with all zeros outside is maximum phase. A linear-phase FIR filter has its
zeros in reciprocal pairs (or on the circle) and sits in between.

The factor that does the reflecting is the first-order \term{all-pass} section
\begin{equation}
A(z)=\frac{z^{-1}-a^*}{1-a\,z^{-1}},\qquad |A(e^{j\omega})|=1,\qquad
\tau_A(\omega)=\frac{1-|a|^2}{|1-a\,e^{-j\omega}|^2}>0,
\label{eq:ch06:allpass}
\end{equation}
which has a pole at $a$ and a zero at $1/a^*$. Its group delay $\tau_A$ is positive at every
frequency and peaks at the angle of the pole. Any rational $H(z)$ factors as $H(z)=H_{\min}(z)
A(z)$, a minimum-phase filter times an all-pass. All-pass filters are the tools for three jobs in
this chapter: equalising the group delay of an IIR filter (Section~\ref{sec:ch06:gdeq}), building
fractional-delay elements, and constructing very efficient half-band and Hilbert filters from
pairs of all-pass branches.

\subsection{FIR or IIR?}

\begin{table}[htbp]\centering\small
\caption{FIR and IIR filters compared.}\label{tab:ch06:firiir}
\begin{tabular}{p{3.2cm}p{5.3cm}p{5.3cm}}
\toprule
& FIR & IIR \\
\midrule
Phase & Exactly linear if coefficients are symmetric & Nonlinear, worst near band edges \\
Stability & Always stable & Must be ensured; sensitive to quantization \\
Order for a sharp filter & High (tens to hundreds of taps) & Low (4 to 10) \\
Multirate efficiency & Excellent (polyphase) & Limited \\
Fixed-point behaviour & Benign & Limit cycles, overflow, scaling \\
Latency & At least half the length for linear phase & Low \\
Typical use in radios & Pulse shaping, channel filters, decimators & DC blockers, AGC and loop filters, audio, notches \\
\bottomrule
\end{tabular}
\end{table}

In communications, the need for linear phase (a filter with nonlinear phase smears pulses in time,
which is intersymbol interference; Chapter~\ref{ch:baseband}) and for efficient rate changes makes
FIR filters the default (Table~\ref{tab:ch06:firiir}). IIR filters earn their place where phase
does not matter, where very low order or very low latency is essential, or where the filter is
inside a feedback loop. A DC blocker, a notch, the loop filter of a PLL (Chapter~\ref{ch:sync})
and the smoothing filter of an AGC are all IIR, and so is the 1-bit $\Sigma\Delta$ modulator's
noise-shaping loop of Chapter~\ref{ch:pcm}.

\begin{history}[How filters became digital]
The digital filter was born as a simulation tool. In the early 1960s engineers at Bell
Laboratories and MIT Lincoln Laboratory wanted to test speech coders---vocoders---without
building every candidate in hardware, so they simulated the analog filter banks on computers.
Doing that faithfully forced them to work out how to turn a continuous-time filter into a
difference equation, and it gradually became clear that the difference equation was a filter in
its own right, with properties no analog circuit could match. James Kaiser at Bell Labs wrote the
first systematic treatment, the chapter on digital filters in Kuo and Kaiser's \emph{System Analysis
by Digital Computer} (1966), including the bilinear transform and a frank discussion of coefficient
sensitivity; his adjustable window, based on the modified Bessel function $I_0$, followed in the
early 1970s and is still the first thing most engineers reach for. Bernard Gold and Charles Rader at
Lincoln Laboratory published \emph{Digital Processing of Signals} in 1969, the first textbook of
the field. When the fast Fourier transform arrived in 1965 (Chapter~\ref{ch:signals}) and integrated
multipliers followed in the 1970s, the simulation tool became the implementation, and the
modem, the vocoder and eventually the whole radio moved into arithmetic.
\end{history}

%======================================================================
\section{Specifying a filter and estimating its cost}
\label{sec:ch06:spec}
%======================================================================

A design starts with a specification, and a surprising amount of engineering judgement goes into
writing a good one. A low-pass specification has four numbers: the passband edge $f_p$, the stopband
edge $f_{st}$, the maximum passband ripple $\delta_p$ (or $A_p$ in dB, peak to peak) and the minimum
stopband attenuation $A_s=-20\log_{10}\delta_s$ in dB. The transition width $\Delta f=f_{st}-f_p$
and $A_s$ largely determine the cost. Converting between the two passband descriptions,
\begin{equation}
\delta_p=\frac{10^{A_p/20}-1}{10^{A_p/20}+1},\qquad A_s=-20\log_{10}\delta_s,
\end{equation}
so $A_p=0.1$\,dB is $\delta_p=0.0058$ and $A_s=60$\,dB is $\delta_s=0.001$.

Where do the numbers come from? In a receiver channel filter, the passband edge is the edge of the
wanted signal (for a root-raised-cosine signal, $(1+\beta)R_s/2$); the stopband edge is where the
nearest interferer begins, or, in a decimator, the frequency that will alias into the band; the
stopband attenuation is set by the strongest adjacent-channel signal the standard requires the
receiver to tolerate plus the SNR it must then achieve; and the passband ripple is set by how much
in-band distortion (EVM) the modem budget allows. A transmit filter is driven instead by the
spectral emission mask and the adjacent-channel leakage ratio (Chapter~\ref{ch:signals}).

\subsection{Order estimates}

Before designing anything, estimate the cost. For an optimal (equiripple) linear-phase FIR
low-pass, three formulas are in common use. fred harris's rule of thumb is the easiest to
remember,
\begin{equation}
N\approx\frac{A_s}{22\,(\Delta f/f_s)},
\label{eq:ch06:harris}
\end{equation}
so a filter with 60\,dB of attenuation and a transition width of 1\% of the sample rate needs
about 270 taps. Bellanger's formula brings in the passband ripple,
\begin{equation}
N\approx\frac{2}{3}\log_{10}\!\left(\frac{1}{10\,\delta_p\delta_s}\right)\frac{f_s}{\Delta f},
\label{eq:ch06:bellanger}
\end{equation}
and Kaiser's formula gives the length of a \emph{Kaiser-window} design,
\begin{equation}
N\approx\frac{A-7.95}{14.36\,\Delta f/f_s}+1,\qquad A=-20\log_{10}\min(\delta_p,\delta_s).
\label{eq:ch06:kaiserN}
\end{equation}
(The more elaborate Herrmann--Rabiner--Chan formula, used by MATLAB's \texttt{firpmord}, is a
refinement of the same idea.) All three agree on the two lessons that drive the rest of this
chapter. First, cost is inversely proportional to the transition width: halving the transition
doubles the length. Second, the transition width is measured \emph{relative to the sample rate at
which the filter runs}. A 10\,kHz transition is 1\% of a 1\,MS/s sample rate but 10\% of a
100\,kS/s rate, where the same filter costs a tenth as many taps, each running ten times less often.
That is the key to the multirate savings of Section~\ref{sec:ch06:multirate}: \emph{run sharp
filters at low rates}.

\begin{keyidea}[Cost is attenuation divided by relative transition width]
The length of an FIR filter grows linearly with the stopband attenuation in dB, only
logarithmically with the passband accuracy, and inversely with the transition width as a fraction
of the sample rate. Asking for 20\,dB more attenuation costs about a third more taps; asking for a
transition twice as sharp costs twice as many. And the multiplication rate is the length times the
rate at which outputs are computed, so the cheapest place for a sharp filter is the lowest sample
rate in the system.
\end{keyidea}

\begin{worked}[Sizing a channel filter]
A 1.4\,MHz LTE carrier occupies 1.08\,MHz (six resource blocks of 180\,kHz), and a receiver
delivers it at 1.92\,MS/s. The channel filter must pass $\pm540$\,kHz with no more than 0.1\,dB of
ripple and attenuate everything beyond $\pm700$\,kHz by 60\,dB. So $\Delta f/f_s=160/1920=0.0833$,
$\delta_p=0.0058$ and $\delta_s=0.001$.
\begin{itemize}
\item harris: $N\approx60/(22\times0.0833)=33$ taps.
\item Bellanger: $\log_{10}\bigl(1/(10\times0.0058\times0.001)\bigr)=4.24$, so
$N\approx\frac23\times4.24/0.0833=34$ taps.
\item Kaiser window: $A=60$\,dB (the stopband is the tighter requirement), so
$N\approx(60-7.95)/(14.36\times0.0833)+1=44.5$, i.e.\ 45 taps with $\beta=0.1102(60-8.7)=5.65$.
\end{itemize}
Running the Parks--McClellan algorithm with the error weighted by $\delta_p/\delta_s=5.8$ and
increasing the length until the specification is met gives 35 taps, as the first two estimates
predicted. The Kaiser design meets the specification with 45 taps, but its passband ripple is only
$\pm0.008$\,dB, far better than required: a window design cannot trade passband accuracy for
length, because the window controls both bands with one parameter
(Figure~\ref{fig:ch06_kaiser_pm}). The 22\% saving matters: at 1.92\,MS/s on complex samples,
35 taps cost $2\times35\times1.92\times10^6=134$ million multiplications per second, or about half
that if the coefficient symmetry is exploited (Section~\ref{sec:ch06:firimpl}).
\end{worked}

%======================================================================
\section{FIR filter design}
\label{sec:ch06:fir}
%======================================================================

\subsection{Linear phase and the four types}

An FIR filter of length $N$ whose coefficients are symmetric, $h[n]=h[N-1-n]$, or antisymmetric,
$h[n]=-h[N-1-n]$, has \term{linear phase}. For odd $N=2L+1$ and symmetric coefficients, pair the
taps around the centre:
\begin{equation}
H(e^{j\omega})=e^{-j\omega L}\Bigl[h[L]+2\sum_{k=1}^{L}h[L-k]\cos(k\omega)\Bigr]
=e^{-j\omega L}A(\omega),
\label{eq:ch06:linphase}
\end{equation}
where the \emph{amplitude} $A(\omega)$ is real (it may be negative, which is a phase jump of $\pi$
that does no harm). Every frequency is delayed by exactly $L=(N-1)/2$ samples, so a pulse passing
through the filter emerges delayed but undistorted. This is the property that makes FIR filters
the default in data transmission, and the reason a root-raised-cosine pair
(Chapter~\ref{ch:baseband}) can be Nyquist end to end. Symmetry also halves the multiplier count:
add $x[n-k]$ and $x[n-N+1+k]$ first, then multiply once.

There are four types, according to whether the length is odd or even and the coefficients
symmetric or antisymmetric (Table~\ref{tab:ch06:types}). The types differ in the zeros they are
\emph{forced} to have. An even-length symmetric filter always has a zero at $z=-1$ ($f_s/2$),
because $A(\omega)$ is then a sum of $\cos\bigl((k-\tfrac12)\omega\bigr)$ terms, all of which vanish at
$\omega=\pi$; it therefore cannot be a high-pass. Antisymmetric filters have $A(\omega)$ built
from sines and are forced to zero at DC, which is exactly right for a differentiator or a Hilbert
transformer, both of which must have zero response at DC, and exactly wrong for a low-pass.

\begin{table}[htbp]\centering\small
\caption{The four types of linear-phase FIR filter. The group delay is $(N-1)/2$ samples in every
case; the antisymmetric types add a constant $90^\circ$ phase shift.}\label{tab:ch06:types}
\begin{tabular}{lllll}
\toprule
Type & Length & Symmetry & Forced zeros & Suitable for \\
\midrule
I & odd & symmetric & none & any: low-, high-, band-pass, half-band \\
II & even & symmetric & $z=-1$ ($f_s/2$) & low-pass, band-pass; not high-pass \\
III & odd & antisymmetric & $z=\pm1$ (DC and $f_s/2$) & Hilbert transformer (band-limited) \\
IV & even & antisymmetric & $z=1$ (DC) & differentiator, Hilbert, high-pass \\
\bottomrule
\end{tabular}
\end{table}

Even-length filters have a group delay of a half-integer number of samples. That is occasionally
useful (an interpolator that must produce samples midway between inputs) and occasionally a trap
(a filter in one branch of a system whose other branch has an integer delay).

\subsection{The window method}

The ideal low-pass filter with cutoff $f_c$ (in cycles/sample) has the infinite impulse response
$h_d[n]=2f_c\sinc(2f_cn)$. Truncating it to $N$ taps is the same as multiplying by a rectangular
window, which convolves the ideal response with the window's spectrum and produces the \emph{Gibbs}
ripple of about 9\% (21\,dB stopband) no matter how long the filter
(Chapter~\ref{ch:signals}). Tapering the truncation with a smooth \term{window} $w[n]$,
\begin{equation}
h[n]=h_d\bigl[n-\tfrac{N-1}{2}\bigr]\,w[n],\qquad
H(e^{j\omega})=\frac{1}{2\pi}\int_{-\pi}^{\pi}H_d(e^{j\theta})\,W\bigl(e^{j(\omega-\theta)}\bigr)\,d\theta,
\end{equation}
trades a wider transition for lower sidelobes: the main lobe of $W$ sets the transition width and
its sidelobes set the stopband ripple. The windows themselves---Hann, Hamming, Blackman,
Blackman--Harris---are the same ones used for spectral analysis in Chapter~\ref{ch:signals}, and
for the same reasons.

The \term{Kaiser window},
\begin{equation}
w[n]=\frac{I_0\Bigl(\beta\sqrt{1-\bigl(\tfrac{2n}{N-1}-1\bigr)^2}\Bigr)}{I_0(\beta)},
\label{eq:ch06:kaiserwin}
\end{equation}
with the shape parameter chosen from the required attenuation $A$ in dB as
\begin{equation}
\beta=\begin{cases}0.1102(A-8.7), & A>50,\\ 0.5842(A-21)^{0.4}+0.07886(A-21), & 21\le A\le 50,\\
0, & A<21,\end{cases}
\label{eq:ch06:kaiser}
\end{equation}
is a near-optimal approximation to the prolate spheroidal sequence, the window that concentrates the
most energy in a given main-lobe width. Its single parameter $\beta$ trades main-lobe width
against sidelobe level continuously, and Kaiser's empirical formulas~\eqref{eq:ch06:kaiser}
and~\eqref{eq:ch06:kaiserN} give $\beta$ and $N$ directly from the attenuation and transition width.
That makes it the quickest route from a specification to a working filter, and it is entirely
adequate for many tasks. It is not optimal: the window controls passband and stopband ripple
together ($\delta_p=\delta_s$, although $\delta_p$ is often far tighter than needed), and the
ripple is largest near the band edge and decays away from it, wasting attenuation where it is not
needed.

\subsection{Frequency sampling}

A second classical method specifies the desired response at $N$ equally spaced frequencies
$\omega_k=2\pi k/N$ and takes the inverse DFT to get the $N$ taps. The resulting filter matches the
specification exactly at those frequencies and interpolates smoothly between them. With samples that
jump from 1 to 0 at the band edge it suffers the same Gibbs ripple as a truncated sinc, but
Rabiner, Gold and McGonegal showed in 1970 that letting one, two or three \emph{transition samples}
take optimised intermediate values raises the stopband attenuation dramatically: one transition
sample gives roughly 40--50\,dB, three give 80--100\,dB. Frequency sampling is the natural method
when the response is specified point by point (an arbitrary equaliser, a measured channel
inverse, the inverse-sinc of a CIC compensator), and \texttt{scipy.signal.firwin2} is a
windowed version of it. It also has a structural interpretation: the filter can be implemented as
a comb $1-z^{-N}$ followed by a bank of resonators, one per nonzero frequency sample, which is
efficient when only a few samples are nonzero and is the conceptual ancestor of the CIC filter.

\subsection{Least-squares design}

The \emph{least-squares} method chooses the coefficients to minimise the weighted integrated
squared error
\begin{equation}
\varepsilon_2=\int_{\text{bands}}W(\omega)\bigl[A(\omega)-D(\omega)\bigr]^2\,d\omega,
\end{equation}
where $D(\omega)$ is the desired amplitude (1 in the passband, 0 in the stopband) and $W(\omega)$ a
weight. Because $A(\omega)$ is linear in the coefficients, setting the gradient to zero gives a set
of linear equations---one matrix solve, no iteration. The result has slightly higher peak error
than the equiripple design of the same length but much less total stopband energy, because its
stopband ripples decay away from the edge. That is the right trade when the stopband contains
broadband noise rather than a single strong interferer at the worst frequency: the noise power
leaking through the stopband is what matters, and that is precisely what least squares minimises.
Least-squares designs (\texttt{scipy.signal.firls}) are also robust for arbitrary multiband shapes
where the equiripple algorithm can fail to converge.

\subsection{Optimal equiripple design: Parks--McClellan}

The \term{Parks--McClellan} algorithm designs the linear-phase FIR filter that minimises the
maximum weighted error,
\begin{equation}
\delta=\min_{h}\;\max_{\omega\in\text{bands}}\;W(\omega)\bigl|A(\omega)-D(\omega)\bigr|.
\end{equation}
The key observation is that the amplitude~\eqref{eq:ch06:linphase} of a type I filter is a
polynomial of degree $L$ in $x=\cos\omega$, because $\cos(k\omega)$ is the Chebyshev polynomial
$T_k(\cos\omega)$. The problem is therefore the classical one of the best uniform (Chebyshev)
approximation of a function by a polynomial, and the \term{alternation theorem} characterises the
answer completely: \emph{a polynomial of degree $L$ is the unique best weighted approximation on a
set of bands if and only if its weighted error attains its maximum magnitude $\delta$, with
alternating signs, at no fewer than $L+2$ frequencies}. The optimum is therefore
\emph{equiripple}: the error oscillates with equal amplitude across each band, and the ripple
heights in the two bands are in the ratio $\delta_p/\delta_s=W_s/W_p$ set by the weights.

The intuition is worth having. Suppose the error peaked at only $L+1$ alternating points. Then one
could find a polynomial of degree $L$ that has the same sign as the error at each of those peaks
(a polynomial of degree $L$ can be made to pass through any $L+1$ values) and subtract a small
multiple of it, lowering every peak at once. Only when there are too many alternations for any
degree-$L$ correction to follow them is the error irreducible.

The \emph{Remez exchange algorithm} turns the theorem into a procedure. Guess $L+2$ extremal
frequencies. Solve the $L+2$ linear equations that force the weighted error to be $\pm\delta$ with
alternating signs at those frequencies (unknowns: the $L+1$ coefficients and $\delta$). Evaluate
the error on a dense grid, find its actual $L+2$ largest alternating peaks, move the guesses
there, and repeat. Each iteration increases $\delta$ toward the minimax value, and convergence
typically takes a handful of iterations. In practice the algorithm is run through barycentric
Lagrange interpolation on a dense grid of about $16N$ points, which is what the 1973 FORTRAN
program by McClellan, Parks and Rabiner did and what \texttt{scipy.signal.remez} still does.

Figure~\ref{fig:ch06_fir_design} compares window, least-squares and equiripple designs of the
same length. At 41 taps the equiripple filter gains about 15\,dB of minimum stopband attenuation
over a Hamming window with a similar transition, and Figure~\ref{fig:ch06_kaiser_pm} shows the
other way of seeing the same advantage: for a fixed specification, the equiripple filter is
about 20\% shorter than the Kaiser-window design.

\mdcfig{ch06_fir_design}{Four 41-tap low-pass designs with a transition band from 0.10 to 0.15
cycles/sample (shaded). The Parks--McClellan filter's stopband is equiripple; the windowed
designs have ripples that decay away from the band edge. Right: the equiripple impulse
response, symmetric and therefore exactly linear phase.}

\mdcfig{ch06_kaiser_pm}{The LTE channel-filter specification of the worked example (pass
$\pm540$\,kHz with $\pm0.05$\,dB ripple, reject beyond $\pm700$\,kHz by 60\,dB, $f_s=1.92$\,MS/s)
met two ways. The Kaiser window needs 45 taps and wastes accuracy: its passband ripple (right) is far
inside the specification and its stopband ripple decays away from the edge. The 35-tap
Parks--McClellan design spends its error budget exactly, touching both limits with equal ripples.}

\begin{history}[Parks, McClellan and an old Russian theorem]
In 1971 Thomas Parks, a young professor at Rice University, set his graduate student James McClellan
the problem of designing optimal linear-phase FIR filters. Earlier approaches had used linear
programming or had fixed the band edges in awkward ways. Parks and McClellan realised that the
problem was exactly the one Evgeny Remez had solved in 1934 for polynomial approximation in the
Chebyshev sense, and that the alternation theorem guaranteed a unique, equiripple optimum. Their
paper ``Chebyshev approximation for nonrecursive digital filters with linear phase'' appeared in
the \emph{IEEE Transactions on Circuit Theory} in March 1972. The following year McClellan, Parks
and Lawrence Rabiner of Bell Labs published the FORTRAN program that implemented it, and it
became one of the most widely copied pieces of software in engineering; its algorithm survives
almost unchanged in \texttt{scipy.signal.remez}, MATLAB's \texttt{firpm} and GNU Radio's
\texttt{optfir}. More than half a century later, ``remez'' is still a verb in DSP labs.
\end{history}

\subsection{Half-band filters}

A \term{half-band filter} is a type I low-pass whose amplitude satisfies
\begin{equation}
A(\omega)+A(\pi-\omega)=1,
\label{eq:ch06:halfband}
\end{equation}
so its response is antisymmetric about the quarter sampling rate $f_s/4$: the transition band is
centred on $f_s/4$, the passband and stopband edges are $f_p$ and $f_s/2-f_p$, and the passband and
stopband ripples are equal. The striking consequence is in the time domain. Index the taps from the
centre, $n=-L,\dots,L$. Replacing $\omega$ by $\pi-\omega$ multiplies $h[n]$ by $(-1)^n$, so the
left side of~\eqref{eq:ch06:halfband} is the transform of $h[n]+(-1)^nh[n]$ (times a linear phase
term), which is $2h[n]$ for even $n$ and zero for odd $n$. For that to equal 1 at every frequency,
\begin{equation}
h[0]=\tfrac12,\qquad h[n]=0\ \text{for even } n\ne0.
\end{equation}
Every second coefficient, except the centre, is exactly zero (Figure~\ref{fig:ch06_halfband}). A
half-band filter of length $N=4K-1$ therefore has only $K$ distinct nonzero non-central
coefficients, and the centre tap is a shift by one bit. Combined with symmetry, a half-band
decimator costs about $N/4$ multiplications per input sample pair---roughly a quarter of a general
filter of the same length.

\mdcfig{ch06_halfband}{A 23-tap equiripple half-band filter. Left: every even-offset coefficient
except the centre is exactly zero (open circles). Centre: the response is antisymmetric about
$f_s/4$, with equal passband and stopband ripple. Right: the zero-phase amplitude and its mirror
image sum to exactly one, which is the defining property~\eqref{eq:ch06:halfband}.}

Half-band filters are designed with Parks--McClellan using symmetric band edges (the algorithm
lands on a half-band solution automatically, to within rounding, and the near-zero taps are then
set exactly to zero), or exactly by designing a filter of half the length for a single band and
interleaving it with zeros. Their natural job is decimation or interpolation by two, where the
image or alias band that must be rejected is the mirror of the passband about $f_s/4$. Cascades of
them, each decimating by two, are the standard way to decimate by powers of two
(Section~\ref{sec:ch06:multistage}).

\subsection{Nyquist (\texorpdfstring{$M$}{M}th-band) filters}

The half-band filter is the $M=2$ member of a family. An \term{$M$th-band} or \term{Nyquist filter}
has
\begin{equation}
h[0]=\frac1M,\qquad h[kM]=0\ \text{for } k\ne0
\quad\Longleftrightarrow\quad
\sum_{k=0}^{M-1}H\bigl(e^{j(\omega-2\pi k/M)}\bigr)=1.
\end{equation}
The time-domain statement is the Nyquist zero-ISI criterion of Chapter~\ref{ch:baseband} in
discrete time: a raised-cosine filter sampled at $M$ samples per symbol is an $M$th-band filter.
In multirate systems, Nyquist filters have two valuable properties. In an interpolator by $M$ the
original samples pass through unchanged (the filter only fills in the missing ones), and one of
the $M$ polyphase branches is a pure delay. And their frequency-domain property---the $M$ shifted
copies sum to a constant---is exactly what a filter bank needs for its channels to tile the band
without gaps (Section~\ref{sec:ch06:filterbanks}).

\subsection{Differentiators and Hilbert transformers}

Two non-low-pass FIR filters appear constantly in radio receivers. A \emph{differentiator} has
$D(\omega)=j\omega$: it computes the derivative of the band-limited signal underlying the samples,
used in FM discriminators (Chapter~\ref{ch:analog}), in the derivative matched filter of
maximum-likelihood timing recovery (Chapter~\ref{ch:sync}) and in edge detectors. A
\emph{Hilbert transformer} has $D(\omega)=-j\operatorname{sgn}\omega$: it shifts every positive
frequency by $-90^\circ$ and is the classical route from a real signal to the analytic signal and
single-sideband modulation (Chapter~\ref{ch:analog}).

Both need antisymmetric coefficients (types III or IV), since both have zero response at DC. A
type IV (even-length) differentiator is best, because it is not forced to zero at $f_s/2$ where
$|D|=\pi$ is largest; its half-sample delay is the price. An odd-length (type III) Hilbert
transformer has a forced zero at $f_s/2$ as well as DC, which is acceptable because the Hilbert
transformer only has to work over a band excluding both. It also has a beautiful connection to the
half-band filter: shift a half-band low-pass up by $f_s/4$ by multiplying its taps by $j^n$. The
result passes positive frequencies and rejects negative ones---a complex filter that produces the
analytic signal directly---and its imaginary part is a type III Hilbert transformer with every
second tap zero. Designing one is designing the other.

\subsection{Design in practice with \texttt{scipy.signal}}

All of the methods above are a few lines of Python. The listing below designs the LTE channel
filter of the worked example three ways and a half-band filter, and shows the two traps that catch
everyone: the frequency normalisation (always pass \texttt{fs} explicitly; without it,
\texttt{firwin} measures frequencies relative to the Nyquist rate but \texttt{remez} relative to
the sample rate in older code), and the weight convention (the ratio of the weights is the
\emph{inverse} of the ratio of the ripples).
\begin{lstlisting}
import numpy as np
from scipy import signal as sg

fs, fp, fst = 1.92e6, 540e3, 700e3          # LTE 1.4 MHz channel filter
dp = (10**(0.1/20) - 1) / (10**(0.1/20) + 1) # 0.1 dB p-p ripple -> 0.0058
ds = 10**(-60/20)                           # 60 dB

# 1. Kaiser window: length and beta from the specification
N, beta = sg.kaiserord(60, (fst - fp) / (fs / 2))   # width relative to Nyquist!
h_kai = sg.firwin(N, (fp + fst) / 2, window=("kaiser", beta), fs=fs)

# 2. Parks-McClellan: weight = inverse ripple ratio
h_pm = sg.remez(35, [0, fp, fst, fs/2], [1, 0], weight=[1, dp/ds], fs=fs)

# 3. Least squares, same length
h_ls = sg.firls(35, [0, fp, fst, fs/2], [1, 1, 0, 0], weight=[1, dp/ds], fs=fs)

# Half-band (23 taps): symmetric band edges, then force the zeros exactly
h_hb = sg.remez(23, [0, 0.2, 0.3, 0.5], [1, 0])
k = np.arange(23) - 11
h_hb[(k % 2 == 0) & (k != 0)] = 0.0

w, H = sg.freqz(h_pm, worN=8192, fs=fs)      # always check the result
print(N, 20*np.log10(abs(H[w >= fst]).max()))
\end{lstlisting}
The last line is the most important. Always plot or measure the designed response against the
specification: the Remez algorithm can fail silently to converge on difficult multiband problems,
and a design that is one tap too short misses the specification by a few decibels with no warning.

\begin{pitfall}[Specifying more than you need]
The most expensive mistake in filter design is not a bad algorithm but a bad specification. A
stopband requirement of 100\,dB when the ADC's own spurious-free dynamic range is 75\,dB, a passband
ripple of 0.01\,dB in front of a demodulator that tolerates 0.5\,dB, a transition band drawn from
the edge of the occupied bandwidth rather than from the edge of the RRC roll-off: each inflates the
filter by tens of percent, and the inflation compounds through every stage of a multirate chain.
Derive every number from the system budget and write down where it came from.
\end{pitfall}

%======================================================================
\section{Implementing FIR filters}
\label{sec:ch06:firimpl}
%======================================================================

\subsection{Direct, transposed and folded structures}

The difference equation $y[n]=\sum_kh[k]x[n-k]$ can be drawn in several ways that compute the same
numbers but behave very differently in hardware (Figure~\ref{fig:ch06_fir_structures}). The
\emph{direct form} is a tapped delay line: the input shifts along a chain of registers, each tap is
multiplied by its coefficient, and an adder tree (or a chain of adders) sums the products. It is
the natural structure for software, where a loop over the taps is a dot product that SIMD
instructions accelerate. Its weakness in hardware is the long combinational path through the adder
chain: every output waits for $N$ additions.

The \emph{transposed form} is obtained by reversing every arrow, swapping inputs and outputs and
replacing branch points with adders (the transposition theorem guarantees the same transfer
function). The input is now broadcast to all multipliers at once, and the delays sit
\emph{between} the adders, so each adder adds one product to a partial sum that arrived from its
neighbour one clock earlier. The critical path is one multiplier and one adder regardless of the
length, so a transposed FIR filter runs at the full clock rate of the device. Its weakness is the
high fan-out of the input.

\begin{figure}[htbp]\centering
\begin{tikzpicture}[
  dl/.style={draw=navy,fill=navy!6,thick,minimum height=0.6cm,minimum width=0.75cm,font=\sffamily\scriptsize},
  mul/.style={draw=navy,thick,circle,inner sep=0.3pt,minimum size=0.48cm,font=\scriptsize},
  add/.style={draw=accent,thick,circle,inner sep=0.3pt,minimum size=0.48cm,font=\scriptsize},
  arr/.style={-{Stealth[length=1.6mm]},thick,navy}]
% ---- direct form
\node[font=\sffamily\small\bfseries,text=navy,anchor=west] at (-0.4,1.3) {(a) Direct form};
\node (x) at (-0.4,0) {\small $x[n]$};
\foreach \i/\xx in {0/0.6,1/2.6,2/4.6,3/6.6} {\coordinate (t\i) at (\xx,0);}
\node[dl] (d1) at (1.6,0) {$z^{-1}$};
\node[dl] (d2) at (3.6,0) {$z^{-1}$};
\node[dl] (d3) at (5.6,0) {$z^{-1}$};
\draw[thick,navy] (x) -- (t0); \draw[arr] (t0) -- (d1); \draw[arr] (d1) -- (d2); \draw[arr] (d2) -- (d3); \draw[thick,navy] (d3) -- (t3);
\foreach \i in {0,1,2,3} {
  \node[mul] (m\i) at ($(t\i)+(0,-0.9)$) {$\times$};
  \node[font=\scriptsize,anchor=east] at ($(m\i)+(-0.25,0)$) {$h_\i$};
  \draw[arr] (t\i) -- (m\i);
  \fill[navy] (t\i) circle (1.3pt);
}
\foreach \i in {1,2,3} {\node[add] (a\i) at ($(t\i)+(0,-1.8)$) {$+$}; \draw[arr] (m\i) -- (a\i);}
\draw[arr] (m0) |- (a1); \draw[arr] (a1) -- (a2); \draw[arr] (a2) -- (a3);
\draw[arr] (a3) -- ++(1.0,0) node[right,font=\small] {$y[n]$};
\node[font=\sffamily\scriptsize,text=accent,anchor=west] at (0.2,-2.45) {critical path: one multiplier and $N-1$ adders};
% ---- transposed form
\begin{scope}[yshift=-4.2cm]
\node[font=\sffamily\small\bfseries,text=navy,anchor=west] at (-0.4,1.1) {(b) Transposed form};
\node (x2) at (-0.4,0) {\small $x[n]$};
\draw[thick,navy] (x2) -- (6.6,0);
\foreach \i/\xx/\hh in {0/0.6/3,1/2.6/2,2/4.6/1,3/6.6/0} {
  \coordinate (u\i) at (\xx,0);
  \node[mul] (n\i) at ($(u\i)+(0,-0.9)$) {$\times$};
  \node[font=\scriptsize,anchor=east] at ($(n\i)+(-0.25,0)$) {$h_\hh$};
  \draw[arr] (u\i) -- (n\i);
  \fill[navy] (u\i) circle (1.3pt);
}
\foreach \i in {1,2,3} {\node[add] (b\i) at ($(u\i)+(0,-1.8)$) {$+$}; \draw[arr] (n\i) -- (b\i);}
\node[dl] (e1) at (1.6,-1.8) {$z^{-1}$};
\node[dl] (e2) at (3.6,-1.8) {$z^{-1}$};
\node[dl] (e3) at (5.6,-1.8) {$z^{-1}$};
\draw[arr] (n0) |- (e1); \draw[arr] (e1) -- (b1); \draw[arr] (b1) -- (e2); \draw[arr] (e2) -- (b2);
\draw[arr] (b2) -- (e3); \draw[arr] (e3) -- (b3);
\draw[arr] (b3) -- ++(1.0,0) node[right,font=\small] {$y[n]$};
\node[font=\sffamily\scriptsize,text=accent,anchor=west] at (0.2,-2.45) {critical path: one multiplier and one adder; input broadcast to all taps};
\end{scope}
% ---- folded (6 taps, 3 multipliers)
\begin{scope}[xshift=9.6cm,yshift=0.3cm]
\node[font=\sffamily\small\bfseries,text=navy,anchor=west] at (-1.0,1.0) {(c) Folded (symmetric, 6 taps)};
\node[font=\small] (xf) at (-0.9,0) {$x[n]$};
\foreach \k/\xx in {0/0,1/1.4,2/2.8} {
  \coordinate (T\k) at ($(\xx,0)+(-0.25,0)$);
  \coordinate (B\k) at ($(\xx,-1.2)+(0.25,0)$);
  \node[add] (P\k) at (\xx,-2.15) {$+$};
  \node[mul] (Q\k) at (\xx,-3.0) {$\times$};
  \node[font=\scriptsize,anchor=east] at ($(Q\k)+(-0.25,0)$) {$h_\k$};
  \draw[arr] (P\k) -- (Q\k);
}
\node[dl] (D1) at (0.7,0) {$z^{-1}$}; \node[dl] (D2) at (2.1,0) {$z^{-1}$};
\node[dl] (D3) at (3.55,-0.6) {$z^{-1}$};
\node[dl] (D4) at (2.1,-1.2) {$z^{-1}$}; \node[dl] (D5) at (0.7,-1.2) {$z^{-1}$};
\draw[thick,navy] (xf) -- (T0); \draw[arr] (T0) -- (D1); \draw[arr] (D1) -- (D2); \draw[thick,navy] (D2) -- (T2);
\draw[arr] (T2) -- (3.55,0) -- (D3); \draw[arr] (D3) -- (3.55,-1.2) -- (D4.east);
\draw[arr] (D4) -- (D5); \draw[thick,navy] (D5.west) -- (B0);
\draw[thick,navy] (D4.east) -- (B2); \draw[thick,navy] (D5.east) -- (B1);
\foreach \k in {0,1,2} {
  \fill[navy] (T\k) circle (1.3pt); \fill[navy] (B\k) circle (1.3pt);
  \draw[arr] (T\k) -- (T\k |- P\k.north west);
  \draw[arr] (B\k) -- (B\k |- P\k.north east);
}
\node[add] (S1) at (1.4,-3.85) {$+$}; \node[add] (S2) at (2.8,-3.85) {$+$};
\draw[arr] (Q0) |- (S1); \draw[arr] (Q1) -- (S1); \draw[arr] (Q2) -- (S2); \draw[arr] (S1) -- (S2);
\draw[arr] (S2) -- ++(0.9,0) node[right,font=\small] {$y[n]$};
\node[font=\sffamily\scriptsize,text=accent,align=left,anchor=north west] at (-1.0,-4.35)
  {pre-add the two samples that share a\\coefficient, multiply once: $N/2$\\multipliers (the DSP-slice pre-adder)};
\end{scope}
\end{tikzpicture}
\caption{FIR filter structures (four taps in (a) and (b), six in (c)). (a) The direct form is a tapped delay line
followed by an adder chain. (b) The transposed form broadcasts the input and pipelines the sum
through registers between the adders, so it runs at the full clock rate however long the filter
is. (c) For a symmetric filter, the folded form adds the two samples that share a coefficient
before multiplying, halving the number of multipliers; FPGA DSP slices include a pre-adder for
exactly this purpose.}
\label{fig:ch06_fir_structures}
\end{figure}

For a symmetric (linear-phase) filter, the \emph{folded} form adds the two delayed samples that
share a coefficient before the multiplication, halving the multiplier count. A half-band filter in
folded form needs only about a quarter of the multipliers of a general filter of the same length.
Hardware designers then decide how much to \emph{parallelise} or \emph{serialise}. If the clock is
much faster than the sample rate, one multiplier can be time-shared across many taps (a
\emph{serial} or \emph{semi-parallel} FIR); if the sample rate exceeds the clock, the filter must
process several samples per clock in parallel, which is how 4\,GS/s RF-sampling converters are
filtered in FPGAs running at 500\,MHz.

\begin{inpractice}[The systolic FIR in an FPGA]
Every modern FPGA contains hundreds to thousands of hard \emph{DSP slices}. A Xilinx DSP48E2 (the
UltraScale family) has a $27\times18$-bit multiplier, a 27-bit pre-adder in front of it and a
48-bit accumulator after it, with pipeline registers at every stage and a dedicated cascade path
that carries a partial sum from one slice to its neighbour without using general routing. A
\emph{systolic} FIR filter is a chain of these slices: each one pre-adds the two symmetric input
samples, multiplies by its coefficient, adds the partial sum arriving on the cascade input and
passes the result on. It is the transposed form with extra pipeline registers distributed along the
input delay line, and it runs at the slice's maximum clock rate---several hundred megahertz---with
no general-fabric logic in the datapath. The 48-bit accumulator means the products can be summed at
full precision and rounded once at the end. Vendor FIR compiler IP generates these structures
automatically, including serialised variants that share one slice among many taps when the clock is
much faster than the sample rate. The older Spartan-6 DSP48A1 used in the USRP B200 has the same
architecture with an $18\times18$ multiplier and an 18-bit pre-adder.
\end{inpractice}

\subsection{Counting the cost}

The cost of a filter is measured in multiplications (or multiply--accumulates, MACs) per second:
taps times the rate at which outputs are computed, times four for a complex filter on complex data
or two for a real filter on complex data.

\begin{worked}[A MAC budget]
A receiver must filter a 20\,MHz LTE carrier sampled at 30.72\,MS/s (complex) with a 127-tap real,
symmetric channel filter. Directly, that is $127\times2\times30.72\times10^6=7.8\times10^9$
multiplications per second (two because I and Q are filtered separately). With folding, the
multipliers halve to 64 per output per rail, $3.9\times10^9$ per second. On an FPGA clocked at
491.52\,MHz (16 times the sample rate), each DSP slice can be time-shared across 16 taps, so the
filter needs $\lceil64/16\rceil\times2=8$ slices. On a CPU core with 256-bit AVX2 vectors (eight
single-precision multiply--adds per fused instruction, two such units per core), the direct form
needs about $10^9$ vector instructions per second, a substantial fraction of one core once the loads
and data shuffling that feed the multipliers are counted, which is why GNU Radio's VOLK library selects the best SIMD kernel for each processor at
run time. Had the same channel filter been applied at a sample rate four times higher, every one of
these numbers would have quadrupled---and the filter would have needed four times as many taps for
the same absolute transition width, so the true cost would have risen sixteen-fold. This is the
argument for decimating before sharp filtering, developed in Section~\ref{sec:ch06:multirate}.
\end{worked}

\subsection{Fast convolution: overlap-add and overlap-save}

For long filters on processors, convolution via the FFT is cheaper. Multiplying DFTs computes
\emph{circular} convolution, so the trick is to process the stream in blocks and arrange that the
circular wrap-around does no harm (Chapter~\ref{ch:signals}). With an $N$-tap filter and an FFT of
size $L\ge N$:
\begin{description}[style=nextline,leftmargin=1.2em,font=\sffamily\bfseries\small\color{navy}]
\item[Overlap-add.] Cut the input into non-overlapping blocks of $L-N+1$ samples, zero-pad each to
$L$, multiply its FFT by the (precomputed) FFT of the zero-padded filter, inverse transform, and
add the overlapping $N-1$-sample tails of consecutive output blocks.
\item[Overlap-save.] Cut the input into blocks of $L$ samples that overlap by $N-1$, transform,
multiply and inverse transform without padding, and discard the first $N-1$ outputs of each block,
which are corrupted by wrap-around. The remaining $L-N+1$ are exactly the linear convolution.
\end{description}
Each block costs two FFTs of size $L$ plus $L$ complex multiplications and produces $L-N+1$
outputs, so the cost per output is about
\begin{equation}
C_{\text{FFT}}\approx\frac{2\cdot\frac{L}{2}\log_2L+L}{L-N+1}\ \text{complex multiplications,}
\end{equation}
minimised for $L$ around four to eight times $N$, giving roughly $2\log_2(4N)$ complex (about
$8\log_2(4N)$ real) operations per output against $N$ real MACs for the direct form. The crossover is
typically between 32 and 128 taps, depending on how well each method vectorises; above a few
hundred taps fast convolution wins by an order of magnitude. Its costs are latency (a whole block
must arrive before any output leaves) and memory. GNU Radio's \texttt{fft\_filter} blocks use
overlap-save. OFDM is, in a sense, fast convolution built into the waveform: the cyclic prefix
makes the channel's linear convolution circular so that one FFT multiplication undoes it
(Chapter~\ref{ch:ofdm}).

%======================================================================
\section{IIR filter design}
\label{sec:ch06:iir}
%======================================================================

\subsection{Analog prototypes}

IIR filters are usually designed by transforming the classical analog prototypes, which were
perfected for LC ladder filters long before digital filters existed. Each family is defined by its
squared magnitude response, normalised to a cutoff of 1\,rad/s:
\begin{description}[style=nextline,leftmargin=1.2em,font=\sffamily\bfseries\small\color{navy}]
\item[Butterworth.] $|H(j\Omega)|^2=1/(1+\Omega^{2n})$. Maximally flat at DC (the first $2n-1$
derivatives vanish), monotonic, and the gentlest roll-off of the family: $-20n$\,dB per decade
asymptotically. Poles equally spaced on a semicircle.
\item[Chebyshev type I.] $|H|^2=1/(1+\varepsilon^2T_n^2(\Omega))$, with $T_n$ the Chebyshev
polynomial. Equiripple in the passband, monotonic in the stopband, considerably sharper than
Butterworth for the same order. Poles on an ellipse.
\item[Chebyshev type II] (inverse Chebyshev). Flat passband, equiripple stopband, with zeros on
the $j\Omega$ axis in the stopband.
\item[Elliptic (Cauer).] Equiripple in both bands, using Jacobi elliptic functions. For a given
order and ripple it has the narrowest transition of any filter---the IIR analogue of the
equiripple FIR---with zeros on the axis that produce deep stopband notches.
\item[Bessel (Thomson).] Maximally flat \emph{group delay} at DC; nearly linear phase in the
passband but a very gentle magnitude roll-off. The bilinear transform spoils its phase property,
so digital designs that need it are usually done directly or replaced by an FIR filter.
\end{description}
The orders needed for a given specification differ dramatically. For a passband ripple of 0.5\,dB,
60\,dB attenuation and a stopband edge 1.5 times the passband edge, a Butterworth filter needs
order 20, a Chebyshev filter order 9 and an elliptic filter order 6. Figure~\ref{fig:ch06_iir_compare}
contrasts three sixth-order designs. The elliptic filter reaches 60\,dB of attenuation in a small
fraction of the transition the Butterworth needs, but its group delay peaks sharply at the band
edge---exactly the kind of phase distortion that causes intersymbol interference in a data
receiver.

\mdcfig{ch06_iir_compare}{Sixth-order IIR low-pass filters with cutoff 0.1 cycles/sample.
Left: magnitude. Centre: group delay, which peaks at the band edge, most severely for the sharpest
design. Right: the elliptic filter's poles cluster near the unit circle and its zeros lie on it.}

\subsection{The bilinear transform}

The \term{bilinear transform} maps an analog transfer function $H_a(s)$ to a digital one by the
substitution
\begin{equation}
s=\frac{2}{T_s}\,\frac{1-z^{-1}}{1+z^{-1}}.
\label{eq:ch06:bilinear}
\end{equation}
It comes from approximating the integrator $1/s$ by the trapezoidal rule, $y[n]=y[n-1]+
\tfrac{T_s}{2}(x[n]+x[n-1])$, whose transfer function is $\tfrac{T_s}{2}(1+z^{-1})/(1-z^{-1})$.
Replace every integrator of an analog filter by its trapezoidal approximation and you have the
bilinear transform. It maps the left half of the $s$-plane onto the inside of the unit circle, so
stable analog filters give stable digital ones, and it maps the entire $j\Omega$ axis onto the unit
circle exactly once, so there is no aliasing: the analog response from $0$ to $\infty$ is squeezed
into $0$ to $f_s/2$. Setting $z=e^{j\omega}$ in~\eqref{eq:ch06:bilinear} gives the price of the
squeezing, a nonlinear \emph{frequency warping},
\begin{equation}
\Omega=\frac{2}{T_s}\tan\frac{\omega}{2}.
\label{eq:ch06:warp}
\end{equation}
At low frequencies $\Omega\approx\omega/T_s$ and the mapping is nearly linear; near $f_s/2$ it
compresses an infinite range of analog frequencies into a finite one. Ripple heights, which are
magnitudes, are preserved exactly; band edges are not. The designer therefore \emph{pre-warps}:
compute the analog band edges from the desired digital ones with~\eqref{eq:ch06:warp}, design the
analog prototype for those edges, and transform. The digital edges then land exactly where they
were wanted.

\begin{worked}[Pre-warping an audio de-emphasis filter]
FM broadcasting in Europe uses a 50\,$\mu$s de-emphasis network, a first-order low-pass
$H_a(s)=1/(1+s\tau)$ with corner $f_c=1/(2\pi\tau)=3183$\,Hz (Chapter~\ref{ch:analog}). A
receiver implements it digitally at $f_s=48$\,kHz. Without pre-warping, the bilinear transform maps
the analog corner $\Omega_c=2\pi\cdot3183$ to the digital frequency
$\omega=2\arctan(\Omega_cT_s/2)=2\arctan(0.2083)=0.4108$\,rad, or 3138\,Hz: a 1.4\% error that
shifts the response by about 0.1\,dB at high audio frequencies. Pre-warping computes
$\Omega_c'=(2/T_s)\tan(\pi\cdot3183/48000)=2\pi\cdot3230$\,rad/s, designs the analog filter for that,
and the digital corner lands exactly at 3183\,Hz. The resulting filter is
\[
H(z)=\frac{0.1745\,(1+z^{-1})}{1-0.6510\,z^{-1}}.
\]
For a corner at a twentieth of the sample rate the error is small; for a filter whose band edges
approach $f_s/4$ or above, ignoring pre-warping produces errors of tens of percent.
\end{worked}

\subsection{Impulse invariance and other mappings}

The alternative classical mapping, \term{impulse invariance}, samples the analog impulse response:
$h[n]=T_sh_a(nT_s)$. Each analog pole $s_k$ becomes a digital pole $e^{s_kT_s}$, and the frequency
axis maps linearly, so the shape of the response is preserved without warping---but the digital
response is the aliased sum of shifted copies of the analog one, so the method works only for
filters that are already well band-limited (low-pass and band-pass, never high-pass). Impulse
invariance is the right tool when the time-domain behaviour matters more than the frequency
response, as in simulating an analog control loop or a resonator. The \emph{matched-$z$} transform
maps poles and zeros alike by $z=e^{sT_s}$; it is simple and useful for quick conversions but
offers no guarantee on the frequency response. Finally, IIR filters can be designed directly in the
digital domain by optimisation (least-squares or minimax fits of a rational function), which is
how filters with simultaneous magnitude and group-delay specifications are usually obtained.

\subsection{Biquad cascades}
\label{sec:ch06:biquads}

An $N$th-order IIR filter could be implemented as one difference equation~\eqref{eq:ch06:diffeq}
(the \emph{direct form}). Mathematically, a cascade of second-order sections---\term{biquads}---is
identical:
\begin{equation}
H(z)=g\prod_{k=1}^{\lceil N/2\rceil}\frac{1+b_{1k}z^{-1}+b_{2k}z^{-2}}{1+a_{1k}z^{-1}+a_{2k}z^{-2}}.
\label{eq:ch06:sos}
\end{equation}
Numerically they are not identical, and the difference is the most important practical fact about
IIR filters. Building the cascade involves three choices:
\begin{enumerate}
\item \textbf{Pairing.} Each complex pole pair is grouped with the zero pair nearest to it, so that
each section's peak gain is moderated by its own zeros. The standard procedure starts with the pole
pair closest to the unit circle (highest $Q$) and pairs it with the nearest zeros, then continues
outwards.
\item \textbf{Ordering.} The sections are ordered by pole radius, conventionally with the
high-$Q$ sections last (to minimise their noise contribution, which the following low-$Q$
sections would otherwise amplify) or first (to minimise the risk of overflow inside the cascade,
because the high-$Q$ sections' peaks are then followed by attenuating ones). Jackson's classic
analysis shows that the choice depends on whether the scaling is by the $L_\infty$ or the $L_2$
norm (Section~\ref{sec:ch06:fixed}); good tools try both.
\item \textbf{Scaling.} Gain is distributed between sections so that no internal node overflows
(Section~\ref{sec:ch06:fixed}).
\end{enumerate}
\texttt{scipy.signal} returns second-order sections directly (\texttt{output="sos"}) and filters
with \texttt{sosfilt}; the transfer-function form \texttt{(b, a)} should be avoided for anything
above fourth order, even in double precision.

\subsection{Numerical behaviour}
\label{sec:ch06:iirnum}

The poles of a high-order polynomial are extremely sensitive to its coefficients when they are
clustered, as they are in any sharp filter. The sensitivity of a root $p_i$ to a coefficient $a_k$
is
\begin{equation}
\frac{\partial p_i}{\partial a_k}=-\frac{p_i^{N-k}}{\prod_{j\ne i}(p_i-p_j)},
\end{equation}
and the product of pole differences in the denominator is tiny when poles are close together.
Figure~\ref{fig:ch06_coef_quant} shows an eighth-order elliptic filter: rounding the direct-form
coefficients to 16 or 12 bits moves poles outside the unit circle and the filter becomes unstable,
while the same filter as four biquads with 12-bit coefficients is almost indistinguishable from the
ideal. \emph{Always implement IIR filters of order above two as cascaded biquads}, with poles paired
to their nearest zeros and sections ordered as above.

\mdcfig[0.72]{ch06_coef_quant}{Coefficient quantization of an eighth-order elliptic filter.
The direct form becomes unstable with 16-bit and 12-bit coefficients; the cascade of biquads
tolerates 12 bits with negligible change.}

Even within a single biquad, the structure matters. In the direct form the coefficients are
$a_1=-2r\cos\theta$ and $a_2=r^2$. Quantizing $a_2$ uniformly quantizes $r^2$, and quantizing $a_1$
quantizes $r\cos\theta$, so the realisable poles lie on the intersections of circles and vertical
lines, which are sparse near $z=1$ (Figure~\ref{fig:ch06_pole_grid}). Narrowband low-pass filters,
whose poles lie exactly there, are the hardest to realise accurately. The \emph{coupled} (Gold--Rader)
form uses the real and imaginary parts of the pole, $r\cos\theta$ and $r\sin\theta$, as coefficients,
giving a uniform grid at the cost of four multipliers instead of two. Lattice and wave digital
structures go further.

\mdcfig{ch06_pole_grid}{Every pole position a second-order section can realise with 5-bit
coefficients. Left: the direct form quantizes $r^2$ and $r\cos\theta$, so the grid is sparse near
$z=1$, exactly where the poles of a narrowband low-pass filter (or a high-$Q$ notch at low frequency)
must go. Right: the coupled form quantizes the real and imaginary parts directly and gives a
uniform grid.}

\begin{description}[style=nextline,leftmargin=1.2em,font=\sffamily\bfseries\small\color{navy}]
\item[Lattice filters.] An all-pole or pole-zero filter can be realised as a cascade of lattice
stages parameterised by \emph{reflection coefficients} $k_m$, the same coefficients that the
Levinson--Durbin recursion produces in linear-predictive speech coding (Chapter~\ref{ch:sourcecoding}).
The filter is stable if and only if every $|k_m|<1$, a condition that is trivially checked and
preserved under quantization, and the structure has excellent sensitivity properties.
\item[Wave digital filters.] Alfred Fettweis's wave digital filters (1971 onwards) simulate a
passive LC ladder using wave variables. Because the prototype is passive, the digital filter
inherits its insensitivity to component values and, with suitable rounding, freedom from limit
cycles and overflow oscillations. They are the most robust IIR structures known, at the price of
design complexity.
\item[All-pass decompositions.] Many odd-order elliptic and Butterworth low-pass filters can be
written as the average of two all-pass filters, $H(z)=\tfrac12[A_0(z)+A_1(z)]$. All-pass sections
are structurally lossless and need only one multiplier per order, so these
\emph{polyphase IIR} or \emph{lattice wave digital} half-band filters are among the cheapest
sharp filters there are, and are used for decimation by two in audio converters.
\end{description}

\subsection{Limit cycles and overflow oscillations}

IIR filters in fixed-point arithmetic exhibit two nonlinear phenomena that FIR filters cannot.
A \term{limit cycle} is a small self-sustained oscillation with zero input, caused by rounding
in the feedback path. Consider a biquad with poles at radius $r=0.98$ whose output has decayed to a
few least-significant bits (LSBs). The ideal output keeps shrinking by 2\% per sample, but when the
products are rounded to integers the decrease of 2\% of a value of 10 LSBs is rounded away: the
effective pole radius becomes 1 and the filter oscillates indefinitely at an amplitude of several
LSBs (Figure~\ref{fig:ch06_limit_cycle}, left). The amplitude of this \emph{granular} limit cycle
is bounded---Jackson's effective-value estimate gives about $0.5/(1-a_2)=0.5/(1-r^2)$ LSBs, or
13 LSBs for $r=0.98$, just what the figure shows---so it matters most in high-$Q$ sections and in applications, such as audio, where a faint
tone in silence is audible. Remedies are magnitude truncation (rounding toward zero, which makes the
effective pole radius smaller rather than larger), extra internal precision, a small dither, or
structures that are provably free of limit cycles.

\mdcfig{ch06_limit_cycle}{Nonlinear behaviour of a fixed-point biquad with zero input. Left: with
products rounded to integers, the decay stalls at a few LSBs and becomes a permanent oscillation
(granular limit cycle). Right: a stable section ($a_1=-1.6$, $a_2=0.9$) started from a large state
oscillates at full scale for ever if its adder wraps around in two's complement, but decays
normally if the adder saturates.}

The second phenomenon is far more dangerous. If an adder inside the feedback loop overflows and
\emph{wraps around}---two's-complement arithmetic turns a large positive number into a large negative
one---the wrap acts as a huge error injected into the loop, and for many coefficient values the
filter then settles into a full-scale \emph{overflow oscillation} that never decays
(Figure~\ref{fig:ch06_limit_cycle}, right). The cure is \emph{saturation arithmetic} in the
recursive part, which clips to the largest representable value instead of wrapping, and is provably
sufficient to prevent overflow oscillations in second-order direct-form sections. Note the contrast
with FIR filters: in a non-recursive sum, intermediate wrap-around is harmless as long as the
final result fits, a property of modular arithmetic that the CIC filter
(Section~\ref{sec:ch06:cic}) exploits deliberately.

\subsection{Group-delay equalisation}
\label{sec:ch06:gdeq}

When a sharp IIR filter is unavoidable---because its latency must be low, or because it is an
analog filter in front of the ADC whose phase must be corrected digitally---its group delay can be
flattened by following it with an all-pass filter whose group delay is the complement of the
filter's. Because $|A(e^{j\omega})|=1$, the magnitude response is untouched, and because every
all-pass section~\eqref{eq:ch06:allpass} adds positive delay, the equalised filter has a larger,
flatter delay than the original. Figure~\ref{fig:ch06_gd_equalize} shows a fifth-order elliptic
filter whose passband group delay varies by about six samples, equalised to a variation of less than
half a sample over 90\% of the passband by three all-pass biquads chosen by numerical optimisation.
The total delay rises from about 6 to about 30 samples, and the cost has risen from three biquads to
six. That is the usual outcome: an elliptic filter plus its delay equaliser is no longer much cheaper
than a linear-phase FIR filter, which is one more reason data receivers usually use FIR filters from
the start.

\mdcfig{ch06_gd_equalize}{Group-delay equalisation. A fifth-order elliptic low-pass (left) has a
group delay that varies by about six samples across the passband and peaks at the band edge (right,
navy). Three all-pass biquads (dotted) add the complementary delay; the total (red) is flat to
within about 0.35 samples over 90\% of the passband, with the magnitude response unchanged. The
price is a larger total delay and twice the arithmetic.}

\begin{pitfall}[IIR filters in the data path]
An elliptic channel filter looks wonderful on a magnitude plot---a sixth-order filter doing the work
of a 60-tap FIR. But a QAM receiver sees the group-delay peak at the band edge as intersymbol
interference that no amount of SNR removes, and an equaliser downstream must spend its taps undoing
it. Before choosing an IIR filter for a digital modem, simulate the EVM it causes, not just its
attenuation. Reserve IIR filters for places where phase does not matter (power detectors, AGC,
spectrum displays), where latency dominates (control loops), or where an all-pass equaliser is
part of the plan.
\end{pitfall}

%======================================================================
\section{Fixed-point arithmetic}
\label{sec:ch06:fixed}
%======================================================================

Floating point is standard on CPUs and GPUs, but FPGAs and ASICs, where most radio DSP at high sample
rates lives, favour fixed point for its cost and power: a fixed-point multiplier is a fraction of
the area and energy of a floating-point one, and an adder is almost free. The price is that the
designer must manage word lengths and scaling explicitly. This section collects what an FPGA or
ASIC designer needs to know.

\subsection{Q formats}

A number in \term{Q-format} $Qm.n$ is a two's-complement integer interpreted with $n$ fractional
bits: $m$ integer bits (including the sign) and $n$ fractional bits, $m+n$ in total. A 16-bit word
in $Q1.15$ represents $[-1,1)$ with resolution $2^{-15}$; a $Q3.13$ word represents $[-4,4)$ with
resolution $2^{-13}$. The rules of the arithmetic follow from integer arithmetic:
\begin{itemize}
\item \textbf{Addition} requires aligned binary points, and the sum of two $Qm.n$ numbers needs
$Q(m{+}1).n$ to be safe from overflow. Adding $K$ numbers needs $\lceil\log_2K\rceil$ extra
integer bits.
\item \textbf{Multiplication} of $Qm_1.n_1$ by $Qm_2.n_2$ gives an exact $Q(m_1{+}m_2).(n_1{+}n_2)$
product. A $16\times16$ product is 32 bits; keeping 16 of them is a quantization.
\item \textbf{Quantization} back to fewer bits is done by \emph{truncation} (dropping the low bits,
which in two's complement rounds toward $-\infty$) or \emph{rounding}.
\end{itemize}

\subsection{Rounding, truncation and quantization noise}

Quantizing to a step $q=2^{-n}$ adds an error $e$ that, for signals large and busy compared with
$q$, behaves like independent uniform noise (Chapter~\ref{ch:pcm}). The three common choices differ
in its mean:
\begin{center}
\begin{tabular}{llp{5.2cm}l}
\toprule
Method & Error range & Mean & Variance \\
\midrule
Truncation (floor) & $(-q,0]$ & $-q/2$ & $q^2/12$ \\
Round half up & $(-q/2,q/2]$ & $\approx0$ (bias $q/2^{k+1}$ when $k$ bits are dropped) & $q^2/12$ \\
Convergent (round half to even) & $[-q/2,q/2]$ & 0 & $q^2/12$ \\
\bottomrule
\end{tabular}
\end{center}
Truncation is free in hardware but adds a DC offset of half an LSB at every quantization point, and
in a receiver that offset appears as a spur at the carrier or at DC after the mixer. Rounding
costs one adder. Convergent rounding removes the residual bias of rounding exact halves upward, which
matters when the same values are rounded millions of times (in an accumulator, or in the long chains
of a decimator); DSP slices implement it with a pattern detector at almost no cost.

In an FIR filter, quantizing each of $N$ products adds $N$ independent noise sources, an output
noise power of $Nq^2/12$. Accumulating at full precision and rounding once at the output adds only
$q^2/12$---a gain of $10\log_{10}N$\,dB, 20\,dB for a 100-tap filter. That is what the wide
accumulators of DSP slices are for. In an IIR filter, the rounding noise injected inside the loop is
shaped by the transfer function from the injection point to the output, so the output noise power
is $\frac{q^2}{12}\sum_n h_e^2[n]$, where $h_e$ is that impulse response; high-$Q$ poles amplify
it by a factor of order $1/(1-r)$, another reason high-$Q$ sections need extra internal bits.

\subsection{Scaling: \texorpdfstring{$L_1$}{L1}, \texorpdfstring{$L_2$}{L2} and \texorpdfstring{$L_\infty$}{L-infinity} norms}

To avoid overflow, the designer chooses a scale factor at each node so that the signal there fits
the word. Three rules are in use, based on three norms of the impulse response $h_k[n]$ from the
input to node $k$:
\begin{equation}
\|h\|_1=\sum_n|h[n]|,\qquad
\|H\|_\infty=\max_\omega|H(e^{j\omega})|,\qquad
\|h\|_2=\Bigl(\sum_nh^2[n]\Bigr)^{1/2},
\end{equation}
which satisfy $\|h\|_2\le\|H\|_\infty\le\|h\|_1$.
\begin{description}[style=nextline,leftmargin=1.2em,font=\sffamily\bfseries\small\color{navy}]
\item[$L_1$ scaling] guarantees no overflow for \emph{any} input bounded by 1, since
$|y[n]|\le\sum|h[k]||x[n-k]|\le\|h\|_1$. It is safe and pessimistic: the worst-case input (the sign
pattern of the time-reversed impulse response) almost never occurs.
\item[$L_\infty$ scaling] guarantees no overflow for any sinusoidal input of unit amplitude. It is
the natural choice for narrowband signals.
\item[$L_2$ scaling] keeps the output \emph{power} of a white input no larger than the input power;
it is a statistical rule that allows occasional overflow (which saturation then handles) in
exchange for better SNR. It suits wideband, noise-like signals---OFDM, for instance.
\end{description}
For a decimation filter in a receiver, a common compromise is to scale for the $L_\infty$ norm of
the passband (the filter's DC gain, usually 1) and add one or two guard bits, relying on saturation
for the rare peaks.

\subsection{Word growth}

The full-precision output of an FIR filter with $B_x$-bit inputs and $B_h$-bit coefficients needs
\begin{equation}
B_y=B_x+B_h+\bigl\lceil\log_2\|h\|_1\bigr\rceil\ \text{bits}
\end{equation}
(with $h$ in integer units), which is at most $B_x+B_h+\lceil\log_2N\rceil$. A 64-tap filter with
16-bit data and 18-bit coefficients needs at most 40 bits---comfortably inside a 48-bit
accumulator. The output is then rounded back to the width the next stage needs, which is set by the
signal-to-noise ratio required there: each bit is worth 6\,dB of dynamic range, and after a
decimation by $M$ that removes out-of-band noise, the in-band SNR has improved by
$10\log_{10}M$\,dB, so the output word can be about $\tfrac12\log_2M$ bits wider than the input
without wasting anything. Coefficient word length is set by the stopband: each coefficient bit buys
roughly 6\,dB of achievable attenuation (Figure~\ref{fig:ch06_coef_bits}), so an 80\,dB filter
needs about 15--16-bit coefficients, and 18-bit coefficients (the native width of most DSP
slices) are good for about 90--100\,dB.

\mdcfig[0.75]{ch06_coef_bits}{Minimum stopband attenuation of a 95-tap equiripple low-pass filter
(designed for about 81\,dB) when its coefficients are rounded to $B$ bits. Each bit is worth about
6\,dB until the design's own attenuation is reached at about 16 bits; beyond that, extra bits buy
nothing.}

\begin{pitfall}[Bit-true simulation]
A filter designed in double precision and implemented in fixed point is a different filter. Before
writing any HDL, simulate the fixed-point implementation \emph{bit-true}---integer arithmetic with
the exact word lengths, rounding modes and saturation of the hardware---and compare its output
against the floating-point model for test vectors that include full-scale tones, full-scale noise
and the worst-case $L_1$ pattern. The comparison catches overflow, bias from truncation and
spurs from coefficient quantization long before they become an FPGA debugging session. NumPy's
integer types (with explicit right shifts), or fixed-point libraries, make this straightforward.
\end{pitfall}

\subsection{A pointer to adaptive filters}

Every filter in this chapter has fixed coefficients. When the coefficients must track an unknown or
changing channel---an equaliser, an echo canceller, an interference canceller, a beamformer---they
are adjusted continuously by an adaptive algorithm such as LMS or RLS. The structures are the same
FIR filters, often in transposed or systolic form, with a coefficient-update loop added; the
fixed-point issues are the same, with the extra hazards of coefficient drift and stalling when the
update is smaller than an LSB. Chapter~\ref{ch:equalization} develops adaptive filtering in full.

%======================================================================
\section{Multirate signal processing}
\label{sec:ch06:multirate}
%======================================================================

A radio's sample rate is chosen at the ADC for the convenience of the analog front end: wide enough
to capture every channel of interest with room for the anti-aliasing filter's transition band. The
demodulator wants something quite different---a few samples per symbol of one channel. Between the
two, the sample rate must change, often by factors of hundreds, and doing that cheaply is the
subject of \emph{multirate} signal processing. Its tools are only two operations, downsampling and
upsampling, plus a handful of identities that let filters be moved past them.

\subsection{Downsampling and aliasing}

To \term{downsample} by an integer $M$ is to keep every $M$th sample, $x_d[m]=x[mM]$. In the
frequency domain,
\begin{equation}
X_d(e^{j\omega})=\frac{1}{M}\sum_{k=0}^{M-1}X\bigl(e^{j(\omega-2\pi k)/M}\bigr).
\label{eq:ch06:downsample}
\end{equation}
The derivation is short: write $x[mM]$ as $x_s[mM]$ where $x_s[n]=x[n]\cdot\frac1M\sum_kW_M^{-kn}$ zeros
every sample except multiples of $M$ (the sum of roots of unity is $M$ at multiples of $M$ and 0
otherwise); multiplying by $W_M^{-kn}=e^{j2\pi kn/M}$ shifts the spectrum by $2\pi k/M$; and
discarding the zeros stretches the frequency axis by $M$. The result is $M$ stretched copies of
the original spectrum, shifted by multiples of $2\pi$. If $X$ is confined to $|\omega|<\pi/M$, the
copies do not overlap and nothing is lost. Otherwise everything between the new Nyquist frequency
$f_s/2M$ and the old one folds into the band, as Chapter~\ref{ch:pcm} warned for the ADC.

A \term{decimator} is therefore a low-pass filter followed by a downsampler
(Figure~\ref{fig:ch06_decimation}). Note what the filter must reject: not everything above
$f_s/2M$, but only what would alias onto the band \emph{that will be used}. If the wanted signal
occupies $|f|<f_p$, the filter's stopband need only begin at $f_s/M-f_p$, because content between
$f_s/2M$ and $f_s/M-f_p$ aliases into the transition region between $f_p$ and $f_s/2M$, where a
later filter can remove it. This ``alias into the transition band'' freedom is what makes half-band
and CIC decimators work.

\mdcfig{ch06_decimation}{Decimation by 4. Downsampling without filtering folds an out-of-band
tone at 0.27 cycles/sample into the band of interest; filtering first removes it.}

\subsection{Upsampling and imaging}

To \term{upsample} by $L$ is to insert $L-1$ zeros between samples, $x_u[n]=x[n/L]$ when $n$ is a
multiple of $L$ and 0 otherwise. Then
\begin{equation}
X_u(e^{j\omega})=X(e^{j\omega L}):
\end{equation}
the spectrum is compressed by $L$, so that the band $|\omega|<\pi$ now contains $L$ copies---the
original at baseband and $L-1$ \emph{images}. An \term{interpolator} is an upsampler followed by a
low-pass filter with cutoff $\pi/L$ that removes the images, with a passband gain of $L$ to restore
the amplitude lost by inserting zeros (a frequent bug: an interpolation filter designed with unit
DC gain produces an output $L$ times too small). A rational rate change by $L/M$ interpolates by
$L$ and then decimates by $M$ with a single filter in between whose cutoff is the smaller of
$\pi/L$ and $\pi/M$.

\subsection{The noble identities}

The two \term{noble identities} let a filter be moved across a rate changer:
\begin{equation}
\underbrace{H(z^M)\ \text{then}\ \downarrow M}_{\text{filter at the high rate}}
\;\equiv\;
\underbrace{\downarrow M\ \text{then}\ H(z)}_{\text{filter at the low rate}},
\qquad
\underbrace{\uparrow L\ \text{then}\ H(z^L)}_{\text{filter at the high rate}}
\;\equiv\;
\underbrace{H(z)\ \text{then}\ \uparrow L}_{\text{filter at the low rate}}.
\label{eq:ch06:noble}
\end{equation}
The first is easy to see in the time domain. $H(z^M)$ is the filter $h$ with $M-1$ zeros inserted
between its taps, so its output at time $mM$ is $\sum_kh[k]x[mM-kM]=\sum_kh[k]x_d[m-k]$, which is
the output of $h$ applied to the downsampled signal. A filter whose taps are spaced $M$ samples
apart only ever sees every $M$th input, so it may as well run at the low rate. The noble
identities apply only to filters that are functions of $z^M$ (or $z^L$); an arbitrary filter
cannot be moved across a rate changer. The trick is to \emph{make} it a function of $z^M$, which
is what polyphase decomposition does.

\subsection{Polyphase decomposition}

Any filter can be split into $M$ interleaved subsequences of its taps,
\begin{equation}
H(z)=\sum_{p=0}^{M-1}z^{-p}E_p(z^M),\qquad e_p[k]=h[kM+p],
\label{eq:ch06:polyphase}
\end{equation}
because the taps $h[p],h[p+M],h[p+2M],\dots$ contribute $z^{-p}\bigl(e_p[0]+e_p[1]z^{-M}+\dots\bigr)$.
The $E_p$ are the \term{polyphase components}. Put~\eqref{eq:ch06:polyphase} in front of a
downsampler: the output is the sum of $M$ branches, branch $p$ being $z^{-p}$, then $E_p(z^M)$,
then $\downarrow M$. By the first noble identity, each $E_p(z^M)$ followed by $\downarrow M$ becomes
$\downarrow M$ followed by $E_p(z)$. The delays and downsamplers together act as a
\emph{commutator} that deals input samples to the branches in turn: branch $p$ receives
$x[mM-p]$. The decimator becomes $M$ short subfilters, each of length $N/M$, each running at the
\emph{output} rate (Figure~\ref{fig:ch06_polyphase_dec}). The total computation is $N$
multiplications per output sample instead of $N$ per input sample---a reduction by a factor of $M$
with no change in the result. The naive decimator computed $M$ outputs and threw $M-1$ away; the
polyphase decimator never computes them.

\begin{figure}[htbp]\centering
\begin{tikzpicture}[
  blk/.style={draw=navy,fill=navy!6,thick,minimum height=0.62cm,minimum width=1.35cm,font=\sffamily\scriptsize},
  add/.style={draw=accent,thick,circle,inner sep=0.3pt,minimum size=0.5cm,font=\scriptsize},
  arr/.style={-{Stealth[length=1.6mm]},thick,navy}]
\node[font=\sffamily\small\bfseries,text=navy,anchor=west] at (-1.2,1.0) {(a) Polyphase decimator};
\node[font=\small] (x) at (-1.0,-1.35) {$x[n]$};
\coordinate (pivot) at (0.15,-1.35);
\foreach \p/\yy/\lab in {0/0/{E_0(z)},1/-0.9/{E_1(z)},2/-1.8/{\vdots},3/-2.7/{E_{M-1}(z)}} {
  \ifnum\p=2
    \node at (2.2,\yy) {$\vdots$};
  \else
    \node[blk] (E\p) at (2.2,\yy) {$\lab$};
    \fill[navy] (0.95,\yy) circle (1.5pt);
    \draw[arr] (1.0,\yy) -- (E\p);
  \fi
}
\draw[thick,navy] (x) -- (pivot); \fill[navy] (pivot) circle (1.5pt);
\draw[-{Stealth[length=1.6mm]},thick,accent] (pivot) -- (0.88,-0.1);
\node[font=\scriptsize,text=accent,align=center] at (-0.35,-2.3) {commutator\\at $f_s$};
\node[add] (S) at (4.0,-1.35) {$+$};
\draw[arr] (E0) -| (S); \draw[arr] (E1.east) -- (S); \draw[arr] (E3) -| (S);
\draw[arr] (S) -- ++(0.9,0) node[right,font=\small] {$y[m]$};
\node[font=\scriptsize,text=navy,align=center] at (2.2,-3.45) {each branch: $N/M$ taps, runs at $f_s/M$};
% ---- interpolator
\begin{scope}[xshift=8.0cm]
\node[font=\sffamily\small\bfseries,text=navy,anchor=west] at (-1.0,1.0) {(b) Polyphase interpolator};
\node[font=\small] (xi) at (-0.8,-1.35) {$x[m]$};
\coordinate (sp) at (0.0,-1.35); \fill[navy] (sp) circle (1.5pt);
\draw[thick,navy] (xi) -- (sp);
\foreach \p/\yy/\lab in {0/0/{E_0(z)},1/-0.9/{E_1(z)},3/-2.7/{E_{L-1}(z)}} {
  \node[blk] (F\p) at (1.4,\yy) {$\lab$};
  \fill[navy] (2.8,\yy) circle (1.5pt);
  \draw[arr] (F\p) -- (2.75,\yy);
}
\node at (1.4,-1.8) {$\vdots$};
\draw[arr] (sp) |- (F0); \draw[arr] (sp) |- (F1); \draw[arr] (sp) |- (F3);
\coordinate (piv2) at (3.65,-1.35); \fill[navy] (piv2) circle (1.5pt);
\draw[-{Stealth[length=1.6mm]},thick,accent] (2.88,-0.1) -- (piv2);
\node[font=\scriptsize,text=accent,align=center] at (4.1,-2.3) {commutator\\at $Lf_s$};
\draw[arr] (piv2) -- ++(0.8,0) node[right,font=\small] {$y[n]$};
\node[font=\scriptsize,text=navy,align=center] at (1.6,-3.45) {branch $p$ computes output phase $p$};
\end{scope}
\end{tikzpicture}
\caption{Polyphase rate changers. (a) A decimator by $M$: a commutator deals successive input samples to
$M$ subfilters $E_p(z)$, whose outputs are summed once per output sample. (b) An interpolator by
$L$: every subfilter sees every input sample, and a commutator collects their outputs in turn, so the
inserted zeros are never multiplied.}
\label{fig:ch06_polyphase_dec}
\end{figure}

The same trick applies to interpolation, using the second noble identity: the $L$ polyphase
subfilters of the interpolation filter each compute one of the $L$ output phases, and the inserted
zeros are never multiplied. This is exactly how a pulse-shaping filter at $L$ samples per symbol is
implemented in a modem: the RRC filter is split into $L$ polyphase branches, and each symbol is
multiplied by each branch once. Polyphase structures are the foundation of efficient DSP, and GNU
Radio's \texttt{pfb\_} (polyphase filter bank) blocks, used in Chapter~\ref{ch:sync} for symbol timing
recovery, are named after them.

\subsection{Rational resampling}

To change the rate by $L/M$, imagine interpolating by $L$ and then decimating by $M$ with one filter
$h$ at the intermediate rate $Lf_s$. Combining the two polyphase tricks, output sample $m$
corresponds to intermediate-rate index $mM$, which lies in input interval $\lfloor mM/L\rfloor$ at
phase $(mM)\bmod L$. So each output is computed by \emph{one} polyphase branch, the one with index
$(mM)\bmod L$, applied to the input samples ending at $\lfloor mM/L\rfloor$. The cost per output is
$N/L$ multiplications, the length of one branch, no matter how large $L$ and $M$ are.

Converting CD audio at 44.1\,kHz to the 48\,kHz of professional and broadcast audio is the classic
case: $48000/44100=160/147$, so $L=160$ and $M=147$. The intermediate rate would be 7.056\,MHz
and the anti-imaging filter, with a transition from 20 to 22.05\,kHz at that rate, needs thousands
of taps. But split into 160 branches of perhaps 30 taps each, the converter costs about 30
multiplications per output sample. \texttt{scipy.signal.resample\_poly} and GNU Radio's
\texttt{rational\_resampler} blocks implement exactly this.

\subsection{Arbitrary and fractional resampling}

Many ratios are irrational or awkward: a 61.44\,MS/s LTE clock to a 56\,MS/s processing rate, a
symbol rate unrelated to the converter clock, or the continuous tracking of a clock that drifts by
parts per million (an \emph{asynchronous} sample-rate converter). Three methods cover these cases.
\begin{itemize}
\item A \emph{polyphase resampler with many phases}---32, 128 or 1024---picks, for each output
sample, the branch nearest to the required fractional delay. With $P$ phases the timing error is at
most $1/2P$ of a sample, which bounds the error at frequency $f$ to roughly $\pi f/P$ relative;
harris's arbitrary resampler (GNU Radio's \texttt{pfb\_arb\_resampler}) adds linear interpolation
between adjacent branches to reach any ratio with 32 phases.
\item A \term{Farrow} structure computes the interpolating filter as a polynomial in the
fractional delay $\mu$, so any delay is synthesised from a few fixed subfilters combined by Horner's
rule. The cubic Lagrange interpolator used for timing recovery in Chapter~\ref{ch:sync} and in
\texttt{commlib} is a four-tap Farrow structure. Its quality depends strongly on oversampling: at
four samples per symbol a cubic interpolator is nearly perfect, at two it is a measurable noise
source.
\item In practice the two are combined: an efficient polyphase stage first interpolates by a
modest factor (say 2 to 4), so that the signal occupies only a small fraction of the band, and a
short Farrow interpolator then does the fine fractional adjustment, where its imperfections are
negligible.
\end{itemize}

\begin{history}[Crochiere, Rabiner, Bellanger and the polyphase idea]
Changing sample rates digitally was well understood in principle by the early 1970s, but doing it
efficiently required a conceptual step. Maurice Bellanger and his colleagues at the French
telecommunications laboratories (TRT) published ``Digital filtering by polyphase network'' in
1976, introducing the polyphase decomposition and showing how it reduces both rate changers and
filter banks for frequency-division-multiplexed telephone channels to banks of short filters and
an FFT---the transmultiplexers that converted between FDM analog groups and TDM digital trunks
(Chapter~\ref{ch:pcm}). At Bell Laboratories, Ronald Crochiere and Lawrence Rabiner developed the
systematic theory of interpolation and decimation, including the optimisation of multistage
designs, in a series of papers culminating in their 1981 \emph{Proceedings of the IEEE} tutorial and
their 1983 book \emph{Multirate Digital Signal Processing}. P.~P.~Vaidyanathan's work on
perfect-reconstruction filter banks in the late 1980s, and his 1993 book, completed the theory.
The ideas spread into every SDR and, through subband coding and the MDCT, into every audio codec
(Chapter~\ref{ch:sourcecoding}).
\end{history}

%======================================================================
\section{Multistage decimation and half-band chains}
\label{sec:ch06:multistage}
%======================================================================

\subsection{Why stages win}

Because FIR cost scales as $A_s/(\Delta f/f_s)$, one large decimation is best done in stages. The
argument is simple. A single-stage decimator by $M$ must have the final, narrow transition band
$\Delta f$ at the full input rate $f_s$, so it needs $N_1\approx K f_s/\Delta f$ taps, where $K$
depends on the ripples; computing one output per $M$ inputs in polyphase form costs
$N_1/M=K f_s/(M\Delta f)$ multiplications per input sample. Now split the decimation into a first
stage by $M_1$ and a second by $M_2=M/M_1$. The first stage only has to protect the band that will
survive the second stage, so its transition can be very wide---from $f_p$ to almost $f_s/M_1$---and
it needs only a handful of taps. The second stage, which must provide the narrow $\Delta f$, runs
at $f_s/M_1$, so it needs $M_1$ times fewer taps for the same absolute transition width, and those
taps are evaluated $M_1$ times less often relative to the input. The sharp filter has moved to the
low rate.

\mdcfig[0.92]{ch06_multistage}{Estimated multiplication rate (per input sample) of a decimate-by-64
filter with 80\,dB of alias protection and a passband to 0.4 of the output sample rate, using
Bellanger's estimate~\eqref{eq:ch06:bellanger} for each stage and polyphase implementation throughout.
A single stage (red) needs about 2100 taps. Splitting into two stages cuts the cost by up to a factor
of five, with the best splits putting the larger factor first; three stages do slightly better, and a
chain of five half-bands followed by a final decimate-by-2 filter (orange) is best of all, and also the
simplest to build.}

Figure~\ref{fig:ch06_multistage} quantifies this for a decimation by 64 with a passband to 0.4
of the output rate and 80\,dB of rejection. A single stage needs about 2100 taps and 33
multiplications per input sample. The best two-stage split ($16\times4$) needs 7; the best
three-stage ones about 6; a chain of half-bands about 5. Two features are typical. First, the
optimum puts the larger factor first, because the first stage's transition is wide anyway and
grows only slowly with its factor, while the last stage's cost is dominated by the narrow final
transition. Crochiere and Rabiner derived closed-form approximations for the optimal factors, but
the cost surface is flat near the optimum, and practical considerations (powers of two, reusable
half-band designs, the clock rate of the hardware) usually decide. Second, the gains saturate after
two or three stages: most of the saving comes from getting the sharp filter off the high-rate
input.

\subsection{Half-band chains}

When the decimation is a power of two, a cascade of half-band decimators is the classic solution,
because each half-band costs roughly a quarter of a general filter and the early ones, with very wide
transitions, are tiny. A half-band decimator in polyphase form has one branch containing only the
centre tap (a delay and a shift) and the other containing the nonzero odd taps, which are
symmetric---so a 4K$-$1-tap half-band costs about $K$ multiplications per \emph{output} sample.

\begin{worked}[A half-band chain from 61.44 to 1.92\,MS/s]
An LTE receiver samples at 61.44\,MS/s (the AD9361's maximum data rate, and a multiple of the
LTE basic rate) but wants the 1.4\,MHz carrier of the earlier worked example at 1.92\,MS/s, a
decimation of $32=2^5$. The wanted signal occupies $\pm540$\,kHz, and we require 80\,dB of
alias protection at each stage. A half-band stage with input rate $F$ must pass $f_p=0.54$\,MHz and
reject from $F/2-f_p$, so its normalised transition is $(F/2-2f_p)/F$. Designing each stage with the
Parks--McClellan algorithm (the shortest length $4K-1$ that achieves 80\,dB) gives:
\begin{center}\small
\begin{tabular}{lcccc}
\toprule
Stage & Rate in $\to$ out (MS/s) & Transition (of $F$) & Taps & Mult./output \\
\midrule
HB1 & $61.44\to30.72$ & 0.482 & 7 & 2 \\
HB2 & $30.72\to15.36$ & 0.465 & 7 & 2 \\
HB3 & $15.36\to7.68$ & 0.430 & 7 & 2 \\
HB4 & $7.68\to3.84$ & 0.359 & 11 & 3 \\
HB5 & $3.84\to1.92$ & 0.219 & 23 & 6 \\
\bottomrule
\end{tabular}
\end{center}
The multiplication rate is $2\times30.72+2\times15.36+2\times7.68+3\times3.84+6\times1.92=131$
million per second per rail---68 per output sample. The 35-tap channel filter of the earlier
example then runs at 1.92\,MS/s and removes what HB5 let alias into the transition region. A
single-stage decimate-by-32 filter with the same 80\,dB rejection and its stopband at
$1.92-0.54=1.38$\,MHz has a relative transition of $0.84/61.44=0.0137$, so~\eqref{eq:ch06:harris}
predicts about 270 taps and about 270 multiplications per output sample in polyphase form: four
times the cost, and one large, awkward filter instead of five small, reusable ones. Notice also that
the first three stages are identical 7-tap filters with two distinct coefficients each; in an ASIC
they can share one design.
\end{worked}

\begin{inpractice}[The AD9361 and the B200's decimation chain]
The AD9361 and AD9364 RF transceivers (the latter in the USRP B200) implement this architecture in
silicon. On the receive side the ADC runs at up to 640\,MS/s, and its output passes through a chain
of fixed decimators---HB3 (which can instead decimate by three, ``DEC3''), HB2 and HB1, each a
half-band decimate-by-two---and then a programmable FIR filter of up to 128 taps that can decimate
by 1, 2 or 4. The transmit side mirrors it: a programmable FIR interpolating by 1, 2 or 4, then HB1,
HB2 and HB3/INT3 into the DAC. The programmable FIR's taps are designed by Analog Devices' filter
wizard for each sample rate, and include compensation for the droop of the analog filters. In the
B200, the FPGA adds a second decimation chain on top: a CORDIC frequency shifter, a CIC decimator and
two half-band filters (Section~\ref{sec:ch06:ddc}). Because the half-bands can be used only when the
FPGA decimation is even, UHD warns when the requested rate implies an odd decimation, and the best
practice is to choose a master clock rate that makes the FPGA decimation a power of two---or one.
\end{inpractice}

%======================================================================
\section{Cascaded integrator-comb filters}
\label{sec:ch06:cic}
%======================================================================

When the decimation ratio is large---hundreds or thousands---even polyphase half-band chains become
expensive, because their early stages run at the full input rate, which may be hundreds of
megahertz. Eugene Hogenauer's \term{cascaded integrator-comb} (CIC) filter (1981) needs no
multipliers and no coefficient storage at all.

\subsection{From moving average to CIC}

Start with the moving average of length $R$, which has the transfer function
\begin{equation}
H_{\text{MA}}(z)=\sum_{k=0}^{R-1}z^{-k}=\frac{1-z^{-R}}{1-z^{-1}}.
\label{eq:ch06:ma}
\end{equation}
The right-hand form is a recursive implementation: an \emph{integrator} $1/(1-z^{-1})$ (an
accumulator) followed by a \emph{comb} $1-z^{-R}$. Its $R$ zeros are evenly spaced around the unit
circle, one of them cancelling the integrator's pole at $z=1$, so the frequency response has nulls
at every multiple of $f_s/R$. Those are precisely the frequencies that alias onto DC after
decimation by $R$. A moving average followed by decimation by $R$ is therefore a natural, if
crude, anti-aliasing decimator: the integrator runs at the input rate, and by the first noble
identity the comb $1-z^{-R}$ followed by $\downarrow R$ is equivalent to $\downarrow R$ followed by
$1-z^{-1}$. The comb now runs at the \emph{output} rate and needs only one delay element.

One stage gives too little attenuation (the first sidelobe of a $\sinc$ is only 13\,dB down), so
Hogenauer cascaded $N$ of them, putting all $N$ integrators at the high rate, then the downsampler,
then all $N$ combs at the low rate (Figure~\ref{fig:ch06_cic_structure}). Allowing a
\emph{differential delay} $D$ (usually 1 or 2) in each comb, the transfer function referred to the
input rate is
\begin{equation}
H(z)=\left(\frac{1-z^{-RD}}{1-z^{-1}}\right)^N,\qquad
|H(e^{j\omega})|=\left|\frac{\sin(\omega RD/2)}{\sin(\omega/2)}\right|^N,
\label{eq:ch06:cic}
\end{equation}
with DC gain $(RD)^N$. Normalised by that gain, the response is approximately $|\sinc(fRD)|^N$
for frequencies small compared with $f_s$: nulls at every multiple of $f_s/RD$, sidelobes falling
as $1/f^N$ (Figure~\ref{fig:ch06_cic}).

\begin{figure}[htbp]\centering
\begin{tikzpicture}[
  blk/.style={draw=navy,fill=navy!6,thick,minimum height=0.7cm,minimum width=0.9cm,font=\sffamily\scriptsize},
  add/.style={draw=accent,thick,circle,inner sep=0.3pt,minimum size=0.46cm,font=\scriptsize},
  arr/.style={-{Stealth[length=1.6mm]},thick,navy}]
\node[font=\small] (x) at (-0.4,0) {$x[n]$};
% integrator 1
\node[add] (a1) at (0.7,0) {$+$};
\node[blk] (d1) at (0.7,-1.05) {$z^{-1}$};
\draw[arr] (x) -- (a1); \draw[arr] (a1.east) -- ++(0.35,0) |- (d1.east); \draw[arr] (d1.north) -- (a1.south);
\node[add] (a2) at (2.2,0) {$+$};
\node[blk] (d2) at (2.2,-1.05) {$z^{-1}$};
\draw[arr] (a1) -- (a2); \draw[arr] (a2.east) -- ++(0.35,0) |- (d2.east); \draw[arr] (d2.north) -- (a2.south);
\node at (3.25,0) {$\cdots$};
\draw[decorate,decoration={brace,amplitude=5pt},navy!70] (0.3,0.45) -- (3.0,0.45)
  node[midway,above=5pt,font=\sffamily\scriptsize,text=navy]{$N$ integrators at $f_s$};
\node[draw=accent,thick,circle,minimum size=0.75cm,font=\scriptsize] (ds) at (4.4,0) {$\downarrow R$};
\draw[arr] (3.55,0) -- (ds);
% comb 1
\node[add] (c1) at (6.4,0) {$+$};
\node[blk] (e1) at (5.6,-1.05) {$z^{-D}$};
\coordinate (b1) at (5.6,0); \fill[navy] (b1) circle (1.4pt);
\draw[arr] (ds) -- (c1); \draw[arr] (b1) -- (e1); \draw[arr] (e1.east) -| node[pos=0.8,right,font=\scriptsize]{$-$} (c1.south);
\node[add] (c2) at (8.2,0) {$+$};
\node[blk] (e2) at (7.4,-1.05) {$z^{-D}$};
\coordinate (b2) at (7.4,0); \fill[navy] (b2) circle (1.4pt);
\draw[arr] (c1) -- (c2); \draw[arr] (b2) -- (e2); \draw[arr] (e2.east) -| node[pos=0.8,right,font=\scriptsize]{$-$} (c2.south);
\node at (9.0,0) {$\cdots$};
\draw[arr] (9.3,0) -- (10.2,0) node[right,font=\small] {$y[m]$};
\draw[decorate,decoration={brace,amplitude=5pt},navy!70] (5.3,0.45) -- (8.9,0.45)
  node[midway,above=5pt,font=\sffamily\scriptsize,text=navy]{$N$ combs at $f_s/R$};
\node[font=\sffamily\scriptsize,text=accent,align=center] at (3.75,-1.25) {registers of\\$B_{\text{in}}+N\log_2(RD)$\\bits};
\end{tikzpicture}
\caption{A CIC decimator: $N$ integrators (accumulators) at the input rate, a downsampler by $R$, and $N$
combs with differential delay $D$ at the output rate. There are no multipliers. Every register must be
wide enough for the full gain $(RD)^N$; wrap-around in the integrators is then harmless.}
\label{fig:ch06_cic_structure}
\end{figure}

\mdcfig[0.72]{ch06_cic}{CIC decimator magnitude response for decimation $R=16$ and one, three
and five stages, plotted against the output sample rate. The nulls sit on the bands that alias to
DC; more stages give deeper nulls and more droop.}

The CIC's strengths are that it needs only $2N$ adders and $N(1+D)$ registers regardless of $R$,
that $R$ can be changed at run time simply by changing the downsampling count, and that the
integrators---the only parts running at the high rate---are the simplest possible circuits. Its
weaknesses are the shape of the response. The attenuation near the nulls is enormous, but the band
that aliases is not a single frequency: when the wanted signal occupies $|f|<f_p$ at the output, the
bands $kf_s/R\pm f_p$ all alias onto it, and the attenuation at the edges of those bands is limited.
For large $R$, the attenuation at the inner edge of the first alias band is about
$\bigl[\sin(\pi b)/\pi(1-b)\bigr]^N$, where $b=f_p/f_{\text{out}}$: with $N=5$, roughly 52\,dB for
$b=\tfrac14$ but 96\,dB for $b=\tfrac1{10}$. That is why a CIC is used for the large, early part of a
decimation, leaving the output still oversampled by four or more, and a conventional FIR filter does
the final, sharp filtering.

\subsection{Droop and its compensation}

The second weakness is \emph{passband droop}: the $\sinc^N$ response falls gradually across the
band of interest. For $R=16$ and $N=4$ the droop at 0.2 of the output rate is 2.3\,dB, far more
than a modem can tolerate. It is corrected by a short FIR \emph{compensation filter} (CFIR) at the
CIC's output rate, designed to have an inverse-$\sinc^N$ passband and usually also to decimate by
two (Figure~\ref{fig:ch06_cic_comp}). Because the compensator runs at the low rate, it is cheap,
and because it can also serve as the first half of the final channel filtering, it is rarely an
additional cost at all. Designing it is a natural job for least squares or frequency sampling
(\texttt{firls} or \texttt{firwin2} with the inverse droop as the desired response), as in the lab.

\mdcfig{ch06_cic_comp}{CIC droop compensation. Left: a CIC ($R=16$, $N=4$) followed by a 31-tap
least-squares compensator that also decimates by two; the cascade is flat to 0.2 of the CIC output rate
and rejects from 0.3 onwards (the band above 0.25, shaded, aliases after the second decimation).
Right: the 2.3\,dB droop is reduced to a ripple of less than 0.07\,dB.}

An alternative is \emph{sharpening}: Kaiser and Hamming showed in 1977 that a filter $H$ with a
nearly flat passband can be improved by the polynomial $3H^2-2H^3$, and applied to a CIC this gives a
\emph{sharpened CIC} with a much flatter passband and deeper alias rejection, at the cost of three
CIC copies and a few shifts and adds---still multiplier-free.

\subsection{Register growth}

The CIC's DC gain is $(RD)^N$, so with $B_{\text{in}}$-bit inputs the output needs
\begin{equation}
B_{\text{out}}=B_{\text{in}}+\bigl\lceil N\log_2(RD)\bigr\rceil
\label{eq:ch06:cicgrowth}
\end{equation}
bits to represent every possible result. The integrators, however, have poles on the unit circle:
their contents grow without bound for any input with a nonzero mean. Hogenauer's key observation was
that this does not matter. Two's-complement arithmetic is arithmetic modulo $2^B$, and the
\emph{overall} CIC is an FIR filter whose output fits in $B_{\text{out}}$ bits. If every register is
$B_{\text{out}}$ bits wide and wraps around on overflow, every intermediate result is correct
modulo $2^{B_{\text{out}}}$, and since the true output fits in that range, the final result is
exactly right. The combs undo the integrators' wrap-arounds. This works only if every stage uses
the same modulus (no saturation anywhere in the CIC) and the registers are at least
$B_{\text{out}}$ bits wide.

Hogenauer also showed that the registers need not all be that wide. Truncating low-order bits at
each stage adds noise, but noise added early is attenuated by the filtering that follows, so the
number of bits that can be discarded (\emph{pruned}) increases from stage to stage. His paper gives
the formulas for choosing the discarded bits so that the total truncation noise at the output is
no more than the output quantization noise; for long decimations the savings are substantial,
often a quarter or more of the register bits.

\begin{worked}[CIC register growth]
A wideband receiver digitises at 245.76\,MS/s with a 14-bit ADC and uses a five-stage CIC ($N=5$,
$D=1$) programmable from $R=8$ to $R=64$ to reach output rates from 30.72\,MS/s down to 3.84\,MS/s.
The register width is set by the largest decimation: $B=14+5\log_2 64=14+30=44$ bits. At $R=8$ the
gain is only $8^5=2^{15}$, so the useful output then lies in the bottom $14+15=29$ bits of the
44-bit word, and the output must be shifted left by $5\log_2(64/8)=15$ bits before rounding to the
output width; for values of $R$ that are not powers of two, $N\log_2R$ is not an integer and the
gain $R^5$ must be normalised by a multiplier (or a shift plus a fine-gain multiplier) after the
combs. Each integrator runs at 245.76\,MHz with a 44-bit adder, which a modern FPGA's carry chains
handle comfortably; the combs run at the output rate and could be time-shared. If the output is
kept to 18 bits, Hogenauer's pruning allows the later integrators and the combs to drop low-order
bits progressively, so that the last comb is only a little wider than the output.
\end{worked}

\subsection{CIC interpolators}

Transposing the CIC gives an interpolator: $N$ combs at the low rate, an upsampler by $R$ (zero
insertion), and $N$ integrators at the high rate. It suppresses the images at multiples of the input
rate with the same $\sinc^N$ shape, and its passband droop is usually pre-compensated by an
inverse-sinc FIR at the low rate. The register-growth analysis is different---the gain from input to
output is $(RD)^N/R$, and the growth happens in the combs---but the wrap-around argument is the same.
CIC interpolators are the last stage of digital up-converters in transmitters, where they raise the
sample rate to the DAC clock.

\begin{history}[Hogenauer's economical filters]
Eugene Hogenauer's paper ``An economical class of digital filters for decimation and interpolation''
appeared in the \emph{IEEE Transactions on Acoustics, Speech, and Signal Processing} in April 1981.
In eight pages it introduced the cascaded integrator-comb structure, derived its frequency response
and alias rejection, proved that integrator overflow is harmless in modular arithmetic, and gave the
register-pruning rules still used today. The timing was perfect. Within a decade, digital receivers
with ADCs sampling at tens of megahertz needed exactly what the CIC offered: a decimator whose
high-rate part was nothing but adders, at a time when a multiplier was an expensive piece of silicon.
The dedicated digital down-converter chips of the 1990s, such as Analog Devices' AD6620 and
Graychip's GC4016, used CIC front ends, and so has almost every DDC since---including the FPGA image
of every USRP. Rarely
has so short a paper been so widely implemented.
\end{history}

%======================================================================
\section{Filter banks and channelizers}
\label{sec:ch06:filterbanks}
%======================================================================

\subsection{The polyphase DFT filter bank}

A bank of $M$ band-pass filters that splits a wideband input into $M$ adjacent channels, each
decimated to its own bandwidth, would naively cost $M$ times one polyphase decimator. The
\term{polyphase DFT filter bank} computes all of them for roughly the cost of one, plus an FFT. The
derivation takes four lines. Let all channels be frequency-shifted copies of one low-pass
\emph{prototype} $h[n]$ with bandwidth about $f_s/M$: channel $k$ uses $h_k[n]=h[n]W_M^{-kn}$, with
$W_M=e^{-j2\pi/M}$, so $H_k(z)=H(zW_M^{k})$. Substituting the polyphase
decomposition~\eqref{eq:ch06:polyphase},
\begin{equation}
H_k(z)=\sum_{p=0}^{M-1}\bigl(zW_M^{k}\bigr)^{-p}E_p\bigl(z^MW_M^{kM}\bigr)
=\sum_{p=0}^{M-1}W_M^{-kp}\;z^{-p}E_p(z^M),
\label{eq:ch06:dftbank}
\end{equation}
because $W_M^{kM}=1$. The polyphase branches $z^{-p}E_p(z^M)$ are \emph{the same for every
channel}; only the complex weights $W_M^{-kp}$ differ, and the sum over $p$ with those weights is an
$M$-point inverse DFT. After the noble identity moves the downsamplers to the front, the whole bank
is a commutator, $M$ polyphase subfilters running at the output rate, and one $M$-point FFT per output
time step (Figure~\ref{fig:ch06_pfb_struct}). The cost per input sample is the prototype length
divided by $M$ plus $\tfrac12\log_2M$ complex multiplications for the FFT, against $M$ times as
much for separate filters. For $M=1024$ channels the saving is nearly a thousandfold.

\begin{figure}[htbp]\centering
\begin{tikzpicture}[
  blk/.style={draw=navy,fill=navy!6,thick,minimum height=0.55cm,minimum width=1.3cm,font=\sffamily\scriptsize},
  arr/.style={-{Stealth[length=1.6mm]},thick,navy}]
\node[font=\small] (x) at (-1.2,-1.2) {$x[n]$};
\coordinate (pv) at (0,-1.2); \fill[navy] (pv) circle (1.5pt); \draw[thick,navy] (x) -- (pv);
\draw[-{Stealth[length=1.6mm]},thick,accent] (pv) -- (0.75,-0.05);
\node[font=\scriptsize,text=accent,align=center] at (-0.5,-2.25) {commutator\\(input rate $f_s$)};
\foreach \p/\yy/\lab in {0/0/{E_0(z)},1/-0.8/{E_1(z)},3/-2.4/{E_{M-1}(z)}} {
  \fill[navy] (0.85,\yy) circle (1.4pt);
  \node[blk] (E\p) at (2.0,\yy) {$\lab$};
  \draw[arr] (0.9,\yy) -- (E\p);
  \draw[arr] (E\p) -- (3.7,\yy);
}
\node at (2.0,-1.55) {$\vdots$};
\node[draw=navy,fill=accent!8,thick,minimum width=1.4cm,minimum height=3.0cm,align=center,font=\sffamily\scriptsize] (F) at (4.4,-1.2) {$M$-point\\IFFT};
\foreach \k/\yy/\lab in {0/0/{channel 0},1/-0.8/{channel 1},3/-2.4/{channel $M{-}1$}} {
  \draw[arr] (5.1,\yy) -- (6.3,\yy) node[right,font=\scriptsize] {\lab};
}
\node at (6.9,-1.55) {$\vdots$};
\draw[decorate,decoration={brace,amplitude=5pt},navy!70] (1.3,0.45) -- (2.7,0.45)
  node[midway,above=5pt,font=\sffamily\scriptsize,text=navy]{prototype split $M$ ways};
\node[font=\sffamily\scriptsize,text=navy,align=center] at (4.4,0.75) {one per $M$ inputs};
\node[font=\sffamily\scriptsize,text=navy,align=center] at (7.0,0.75) {each at $f_s/M$};
\end{tikzpicture}
\caption{A critically sampled polyphase DFT channelizer. One commutator, $M$ short subfilters (the
polyphase components of a single prototype low-pass) and one FFT per output time step produce all $M$
channels, each decimated by $M$ and shifted to baseband.}
\label{fig:ch06_pfb_struct}
\end{figure}

Figure~\ref{fig:ch06_polyphase} shows the channel responses of an 8-channel bank, and
Figure~\ref{fig:ch06_channelizer} shows one at work: a wideband input carrying five carriers of
different powers on an 8-channel grid is split into eight baseband outputs, each at an eighth of the
input rate. (The figure script checks the channelizer against brute-force mixing, filtering and
downsampling of each channel; they agree to within rounding error.)

\mdcfig[0.72]{ch06_polyphase}{An 8-channel uniform filter bank: every channel is a frequency-shifted
copy of one prototype low-pass filter, and a polyphase-plus-FFT structure computes them all at once.}

\mdcfig{ch06_channelizer}{A polyphase DFT channelizer at work. Top: a wideband input with five
QPSK carriers of different power on an 8-channel grid (channels 0, 4 and 7 are empty). Bottom: the
eight outputs, each brought to baseband and decimated by 8 by one polyphase filter bank and one
8-point FFT per output sample. Because this bank is critically sampled, the carriers' roll-off
regions fold back at the edges of each output band.}

\subsection{Critically sampled and oversampled banks}

The bank of Figure~\ref{fig:ch06_pfb_struct} is \emph{critically sampled}: $M$ channels, each
decimated by $M$, so the total output rate equals the input rate. Its weakness is the channel
edges. The prototype's transition band must lie somewhere, and in a critically sampled bank
whatever lies in it aliases into the neighbouring channel's output (visible at the band edges in
Figure~\ref{fig:ch06_channelizer}). A signal that straddles two channels is damaged in both.
\emph{Oversampled} banks avoid this by decimating each channel by $M/2$ (or $M/4$) instead of $M$:
the output rate per channel is then twice the channel spacing, the transition band can fall in the
spare spectrum, and a signal anywhere in the band can be recovered, if necessary by recombining two
adjacent channels with a small synthesis bank. The price is a doubled output rate and a slightly
more complex commutator, which must also apply a circular shift to the FFT input on alternate
steps. harris's book gives the full recipe. Channel \emph{stacking} (whether the channel centres
fall on $kf_s/M$ or on $(k+\tfrac12)f_s/M$) is the other design choice; the latter, implemented
with a pre-rotation by $W_{2M}^{n}$, avoids a channel straddling DC.

\subsection{Synthesis and perfect reconstruction}

The transpose of an analysis bank is a \emph{synthesis} bank: an IFFT, polyphase subfilters and an
output commutator that combine $M$ low-rate channels into one wideband signal. A transmitter that
must generate many carriers at once---a multicarrier base station, a cable headend, a satellite
gateway---uses one. An analysis bank followed by a synthesis bank is a \emph{transmultiplexer} in
reverse: it splits a signal into subbands and reassembles it, and with suitably designed prototypes
the output can be a perfect delayed copy of the input. The theory of such \emph{perfect
reconstruction} filter banks---quadrature mirror filters, paraunitary banks, cosine-modulated
banks---was developed in the 1980s by Smith and Barnwell, Vaidyanathan, Vetterli and others. It
underpins subband audio coding and the MDCT of every modern audio codec (Chapter~\ref{ch:sourcecoding}),
the wavelet transform, and the filter-bank multicarrier (FBMC) waveforms studied as alternatives to
OFDM (Chapter~\ref{ch:ofdm}).

\begin{inpractice}[Channelizers in orbit, in telescopes and in base stations]
Polyphase channelizers sit wherever many narrowband signals share one wide band. Satellite
\emph{digital transparent processors} (Chapter~\ref{ch:satellite}) digitise each uplink beam,
channelize it into many sub-channels, route each sub-channel to any downlink beam and frequency, and
resynthesise the downlink with a synthesis bank; the Inmarsat-4 satellites launched from 2005 were
early large-scale examples, and modern high-throughput and software-defined payloads have taken
the idea much further. Spectrum-monitoring and signals-intelligence receivers channelize an entire
band at once to detect and record every transmission in it. Radio telescopes use polyphase filter
banks as their spectrometers and as the first stage of their correlators; the open-source CASPER
FPGA libraries made polyphase filter-bank spectrometers standard equipment at observatories around
the world. And a software base station that serves many narrowband carriers---GSM, or narrowband IoT
channels---from one wideband ADC uses a channelizer instead of one DDC per carrier.
\end{inpractice}

\subsection{Spectral analysis: Welch and the polyphase FFT}

The DFT itself is a filter bank: each bin is the output of a filter whose impulse response is the
window, modulated to the bin frequency and decimated by the FFT length (Chapter~\ref{ch:signals}).
With a window of length $M$---one FFT block---the bin filter has the window's transform as its
response: a main lobe one to four bins wide and sidelobes that leak strong signals into distant bins.
\emph{Welch's method} averages the squared magnitudes of many overlapping windowed FFTs to reduce the
variance of the estimate, but does nothing for the shape of each bin's response.

A \emph{polyphase FFT}, or polyphase filter-bank spectrometer, uses a prototype $T$ times longer than
the FFT---typically $T=4$ to 16 taps per branch---and folds it onto the FFT with the polyphase
structure of Figure~\ref{fig:ch06_pfb_struct}. Each bin then has the response of a properly designed
low-pass filter: flat across the bin, with steep skirts and low sidelobes
(Figure~\ref{fig:ch06_pfb_spectrum}). Scalloping loss nearly disappears, and leakage from a strong
signal into distant bins drops by tens of decibels, for an extra cost of only $T$ multiplications per
input sample. For the price of a slightly more complicated front end, every bin becomes a good
channel filter, which is why radio astronomers and spectrum monitors prefer the polyphase FFT to the
windowed FFT.

\mdcfig[0.82]{ch06_pfb_spectrum}{The response of one spectrometer channel as a function of frequency
offset from its centre. A rectangular-window FFT has $-13$\,dB sidelobes and 4\,dB of scalloping at the
bin edges (dotted); a Hann window lowers the sidelobes but widens the main lobe. A polyphase FFT with
four taps per branch (a 256-tap windowed-sinc prototype for a 64-point FFT) gives a flat-topped bin with
steep skirts and sidelobes below $-75$\,dB.}

%======================================================================
\section{Numerically controlled oscillators}
\label{sec:ch06:nco}
%======================================================================

\subsection{The phase accumulator}

Digital mixing multiplies samples by $e^{-j2\pi f_0n/f_s}$. A \term{numerically controlled
oscillator} (NCO), or direct digital synthesiser (DDS) when its output drives a DAC, generates that
complex exponential with three components: a $B$-bit \emph{phase accumulator} that adds a
\emph{frequency tuning word} (FTW) $F$ every clock, a \emph{phase-to-amplitude converter} (a sine
look-up table or a CORDIC), and, in a DDS, a DAC. The accumulator wraps around modulo $2^B$, which is
exactly the modulo-$2\pi$ behaviour of phase, so the output frequency and its resolution are
\begin{equation}
f_0=\frac{F}{2^B}\,f_s,\qquad \Delta f=\frac{f_s}{2^B}.
\label{eq:ch06:nco}
\end{equation}
A 32-bit accumulator clocked at 245.76\,MHz has a resolution of 0.057\,Hz; a 48-bit accumulator, used
in precision synthesisers, reaches micro-hertz. Frequency changes are \emph{phase-continuous}---the
accumulator simply starts adding a different number---and take effect within one clock, which is
why NCOs generate the frequency hops of military radios and Bluetooth, the chirps of FMCW radars and
the carriers of many transmitters, and why every DDC uses one as its local oscillator. A frequency
that is not an integer multiple of $\Delta f$ cannot be represented exactly; a 1\,MHz tone from a
245.76\,MHz clock has a fractional tuning word, and the residual error ($\le\Delta f/2$) accumulates
as a slow phase drift that matters in coherent systems. Some DDS chips offer a programmable modulus to
make such frequencies exact.

\subsection{The look-up table and phase truncation}

The phase-to-amplitude converter cannot use all $B$ accumulator bits: a $2^{32}$-entry table is
absurd. Only the top $P$ bits address the table, and the discarded $B-P$ bits are a \emph{phase
truncation} error. For most tuning words that error is a periodic sawtooth (the discarded bits form
a periodic sequence), and a periodic phase error produces \emph{spurs}---discrete spectral lines---
rather than noise. Their worst-case level relative to the carrier is approximately
\begin{equation}
\text{SFDR}\approx6.02\,P\ \text{dBc}
\label{eq:ch06:ncosfdr}
\end{equation}
(Figure~\ref{fig:ch06_nco_spurs}), and the actual level and the spur positions depend on the tuning
word: if the low $B-P$ bits of $F$ are zero there is no truncation error at all, while words whose
low bits are a small fraction of $2^{B-P}$ produce the largest spurs. Amplitude quantization of the
table entries to $A$ bits adds its own error, which behaves like noise at about $6.02A+1.76$\,dB below
full scale for most tuning words but can also concentrate into harmonics when the output period is
short.

\mdcfig{ch06_nco_spurs}{Output spectra of an NCO with a 32-bit accumulator when the phase is
truncated to 8 and to 14 bits before the sine look-up. Each additional phase bit lowers the spurs
by about 6\,dB.}

Three techniques reduce the cost of a large table. \emph{Quarter-wave symmetry} stores only the first
quadrant of the sine and derives the others by inverting the address and the sign, cutting the table
by four. \emph{Sine--cosine sharing} reads both outputs from the same table offset by a quarter
period. And \emph{interpolation} (Taylor-series correction) uses the discarded phase bits with a
small multiplier to refine the table output, so that a table of $2^{10}$ entries can deliver the
accuracy of one of $2^{16}$. Alternatively, a CORDIC (Section~\ref{sec:ch06:cordic}) replaces the
table altogether.

\subsection{Phase dithering}

If the spurs, rather than the total error power, are the problem---as they are in a receiver where a
spur can mask a weak signal or in a transmitter with a spectral mask---the periodicity of the
truncation error can be broken by adding a small random number to the accumulator output before
truncation. The \emph{phase dither}, uniform over one LSB of the truncated phase, turns the
periodic error into white noise: the total error power rises slightly but is spread across the whole
band, and the spurs drop into a noise floor (Figure~\ref{fig:ch06_nco_dither}). With 10 phase bits,
dithering trades $-60$\,dBc spurs for a floor some 25\,dB lower in each FFT bin. The same idea, in the
form of noise shaping, is used in fractional-$N$ synthesisers (Chapter~\ref{ch:sync}).

\mdcfig{ch06_nco_dither}{Phase dithering. Left: with the phase truncated to 10 bits, the error is
periodic and produces spurs up to about $-60$\,dBc. Right: adding a random value below one phase LSB
before truncation converts the spurs into a noise floor; the worst remaining peak is a noise peak
about 25\,dB lower.}

\begin{worked}[An NCO for a 75\,dB receiver]
A digital receiver needs a local oscillator whose spurs are at least 70\,dB below the carrier,
running at 245.76\,MS/s with 0.1\,Hz resolution. Resolution: $f_s/2^B\le0.1$\,Hz requires
$2^B\ge2.46\times10^9$, so $B=32$ (0.057\,Hz). Spurs: by~\eqref{eq:ch06:ncosfdr}, $P=12$ phase bits
give about $6.02\times12=72$\,dBc, just meeting the requirement. The table then has $2^{12}=4096$
entries; with quarter-wave symmetry, 1024 entries of the first quadrant. For the amplitude, $A=14$
bits gives a quantization floor near $6.02\times14+1.76=86$\,dB below full scale, comfortably below
the phase spurs. The table is $1024\times13$ bits (the sign is handled by the symmetry logic), about
13\,kbit, which fits in a single 18\,kbit FPGA block RAM; with sine and cosine read from two ports
of the same memory, the complete quadrature NCO costs one block RAM and a 32-bit adder. If the system
later needs 90\,dBc, three more phase bits would multiply the table by eight, and phase dithering
or Taylor correction becomes the cheaper option.
\end{worked}

\begin{inpractice}[DDS chips and NCOs everywhere]
Stand-alone DDS chips put the whole structure---a 32- or 48-bit accumulator, a phase-to-amplitude
converter and a high-speed DAC---on one die. Analog Devices' AD9910, for example, clocks at up to
1\,GS/s with a 32-bit tuning word and a 14-bit DAC, and can sweep, hop and modulate under external
control; such parts generate the clocks of test equipment, the chirps of laboratory radars and
the agile local oscillators of frequency-hopping radios. Inside an RF-sampling data converter, such as
the ADCs and DACs of Xilinx RFSoC devices or the AD9081, NCOs with 48-bit accumulators and phase
dithering tune each digital down- or up-converter. In software, GNU Radio's \texttt{sig\_source}
and \texttt{rotator} blocks are NCOs whose phase is a floating-point number---a reminder that a
floating-point phase accumulator loses resolution as its value grows, which is why careful software
keeps the phase as an integer or wraps it every sample.
\end{inpractice}

%======================================================================
\section{CORDIC}
\label{sec:ch06:cordic}
%======================================================================

\subsection{Rotation by shifts and adds}

Rotating a vector $(x,y)$ by an angle $\theta$,
\[
\begin{pmatrix}x'\\y'\end{pmatrix}=\cos\theta\begin{pmatrix}1&-\tan\theta\\ \tan\theta&1\end{pmatrix}
\begin{pmatrix}x\\y\end{pmatrix},
\]
needs four multiplications and the sine and cosine of $\theta$. Jack Volder's \term{CORDIC} (COordinate
Rotation DIgital Computer) algorithm does it with shifts and adds by rotating through a fixed
sequence of \emph{elementary angles} $\theta_i=\arctan2^{-i}$, $i=0,1,2,\dots$, each in the
direction $d_i=\pm1$ that brings the result closer to the target. Because $\tan\theta_i=2^{-i}$, each
elementary rotation is
\begin{equation}
x_{i+1}=x_i-d_i\,2^{-i}y_i,\qquad y_{i+1}=y_i+d_i\,2^{-i}x_i,\qquad z_{i+1}=z_i-d_i\theta_i,
\label{eq:ch06:cordic}
\end{equation}
where multiplication by $2^{-i}$ is an arithmetic right shift and $z_i$ tracks the remaining angle
(the $\theta_i$ are a small table of constants). Omitting the factor $\cos\theta_i$ from each step
makes every iteration stretch the vector by $1/\cos\theta_i=\sqrt{1+2^{-2i}}$, independent of the
direction chosen, so after $n$ iterations the total gain is a constant,
\begin{equation}
\frac1K=\prod_{i=0}^{n-1}\sqrt{1+2^{-2i}}\;\longrightarrow\;1.64676,\qquad K\approx0.60725,
\end{equation}
which is removed by one multiplication at the end, or folded into the initial value or a later
gain stage.

The algorithm has two modes:
\begin{description}[style=nextline,leftmargin=1.2em,font=\sffamily\bfseries\small\color{navy}]
\item[Rotation mode] chooses $d_i=\operatorname{sgn}(z_i)$, driving the residual angle $z$ to zero.
Starting from $(x_0,y_0,z_0=\theta)$, it ends with the input vector rotated by $\theta$ (and scaled by
$1/K$). Starting from $(K,0,\theta)$ it produces $\cos\theta$ and $\sin\theta$: an NCO without a
table. Applied to each sample with $\theta$ from a phase accumulator, it is a complete digital mixer.
\item[Vectoring mode] chooses $d_i=-\operatorname{sgn}(y_i)$, driving $y$ to zero: it rotates the
vector onto the $x$ axis. It ends with $x_n=\sqrt{x_0^2+y_0^2}/K$ and $z_n=z_0+\arctan(y_0/x_0)$,
computing magnitude and phase in one pass. That is exactly what AM and FM demodulators
(Chapter~\ref{ch:analog}), AGC loops and the phase detectors of synchronisers (Chapter~\ref{ch:sync})
need.
\end{description}
The elementary angles sum to $\sum_i\arctan2^{-i}\approx1.7433$\,rad (about $99.9^\circ$), so the
basic algorithm converges for $|\theta|<99.9^\circ$; a preliminary rotation by $\pm90^\circ$ (swap
$x$ and $y$ and negate one) extends it to the full circle. Generalisations by John Walther of
Hewlett-Packard (1971) use the same iteration in linear and hyperbolic coordinates to compute
multiplication, division, exponentials, logarithms and square roots.

\subsection{Accuracy and cost}

Each iteration halves the size of the elementary rotation, so the residual angle after $n$
iterations is bounded by $\theta_{n-1}=\arctan2^{-(n-1)}\approx2^{-(n-1)}$\,rad: CORDIC delivers
roughly \emph{one bit of accuracy per iteration} (Figure~\ref{fig:ch06_cordic}). The datapath must be
wider than the desired accuracy, because each iteration's right shift discards bits and the rounding
errors of $n$ iterations accumulate; the usual rule is $\log_2n+2$ guard bits. In hardware the
iterations are unrolled into a pipeline of $n$ identical stages, each with three adders and fixed
shifts, producing one result per clock with a latency of $n$ clocks; or a single stage is reused
$n$ times when throughput is low.

\mdcfig{ch06_cordic}{CORDIC. Left: rotation-mode trajectory for a target of $57^\circ$, starting from
$(1,0)$; each step rotates by $\pm\arctan2^{-i}$ and stretches the vector, so the path spirals out to
radius $1/K\approx1.647$. Right: maximum phase error over 4000 random angles as a function of the
number of iterations, for 24-bit and 18-bit integer datapaths. The error halves with each iteration
until the datapath's resolution is reached.}

\begin{worked}[How many iterations for 16-bit accuracy?]
A DDC needs a CORDIC mixer whose phase error is below one LSB of a 16-bit output, about
$2^{-15}=3.1\times10^{-5}$\,rad. The residual angle after $n$ iterations is at most
$\arctan2^{-(n-1)}$, so $n-1\ge15$, i.e.\ $n=16$ iterations. With the guard-bit rule, the internal
datapath needs about $16+\log_216+2=22$ bits. The simulation behind Figure~\ref{fig:ch06_cordic}
confirms it: with a 24-bit datapath, 16 iterations give a maximum error of $3.2\times10^{-5}$\,rad,
while an 18-bit datapath stalls near $1.5\times10^{-4}$\,rad no matter how many iterations are run.
The pipelined implementation is 16 stages of three 22-bit adders---48 adders and no multipliers, one
output per clock at several hundred megahertz in an FPGA---plus a final multiplication by
$K=0.60725$ if the gain must be exactly one (or none, if the following stage can absorb it). A
look-up-table NCO with a complex multiplier needs about the same logic, which is why both
appear in practice; UHD's FPGA DDC uses a CORDIC.
\end{worked}

\begin{history}[Volder's navigation computer]
In the mid-1950s Jack Volder, an engineer at Convair's aeronautics division in Fort Worth, was asked to
replace the analog resolver in the navigation system of the B-58 Hustler, the first supersonic
bomber, with a digital computer. Navigation by dead reckoning required continuous conversion between
polar and rectangular coordinates---headings and distances to north and east components and back---and
the transistor computers of the day had no fast multiplier. Volder found an iteration, which he related to an
old method for computing trigonometric tables, that needed only shifts and additions, and described it in
``The CORDIC trigonometric computing technique'' in the \emph{IRE Transactions on Electronic Computers}
in September 1959. John Walther's 1971 unification extended it to the full range of elementary
functions, and shift-and-add methods of this family computed the transcendental functions of early
scientific pocket calculators. Today CORDIC lives on wherever multipliers are scarce or rotations are
the natural operation: in FPGA IP libraries, in the digital front ends of radio chips, and in
the phase detectors of countless synchronisers.
\end{history}

%======================================================================
\section{Digital down- and up-converters}
\label{sec:ch06:ddc}
%======================================================================

\subsection{The DDC architecture}

Everything in this chapter comes together in the \term{digital down-converter} (DDC), the block that
takes a wideband stream of samples from an ADC and delivers one channel at baseband at a low sample
rate (Figure~\ref{fig:ch06_ddc}). Its stages are now familiar:
\begin{enumerate}
\item A \textbf{complex mixer}, driven by an NCO or a CORDIC, shifts the wanted channel to 0\,Hz. If the
ADC samples a real IF, the mixer produces I and Q from one real input; if the front end is already
zero-IF, the mixer fine-tunes by a few megahertz.
\item A \textbf{CIC decimator} provides the bulk of the decimation, from the ADC rate down to four or
eight times the final rate, with no multipliers.
\item A \textbf{compensating FIR} (CFIR) corrects the CIC droop and decimates by two.
\item A \textbf{programmable channel FIR} (PFIR) sets the final channel shape---often a root-raised
cosine matched filter---and decimates by two again.
\item Optionally, a \textbf{fractional resampler} converts to an exact symbol-related rate, and an
\textbf{AGC} scales the output word.
\end{enumerate}

\begin{figure}[htbp]\centering
\begin{tikzpicture}[node distance=0.32cm,
  blk/.style={draw=navy,fill=navy!6,thick,minimum height=1.0cm,minimum width=1.25cm,align=center,font=\sffamily\scriptsize},
  arr/.style={-{Stealth[length=1.8mm]},thick,navy}]
\node[blk,fill=accent!6,draw=accent] (adc) {ADC\\14 bit};
\node[draw=navy,thick,circle,minimum size=0.6cm,right=0.5cm of adc,font=\small] (mix) {$\times$};
\node[blk,below=0.55cm of mix] (nco) {NCO /\\CORDIC};
\node[blk,right=0.5cm of mix] (cic) {CIC\\$\downarrow R$\\($N=5$)};
\node[blk,right=of cic] (cfir) {CFIR\\$\downarrow2$\\droop corr.};
\node[blk,right=of cfir] (pfir) {PFIR\\$\downarrow2$\\channel/RRC};
\node[blk,right=of pfir] (rs) {resampler\\$\times L/M$\\(Farrow)};
\node[blk,right=of rs] (agc) {AGC};
\draw[arr] (adc) -- (mix); \draw[arr] (nco) -- (mix); \draw[arr] (mix) -- node[above,font=\scriptsize]{I/Q} (cic);
\draw[arr] (cic) -- (cfir); \draw[arr] (cfir) -- (pfir); \draw[arr] (pfir) -- (rs); \draw[arr] (rs) -- (agc);
\draw[arr] (agc.east) -- ++(0.5,0) node[right,font=\scriptsize,align=left]{baseband\\out};
\node[font=\sffamily\scriptsize,text=accent,below=0.12cm of adc] {$f_s$};
\node[font=\sffamily\scriptsize,text=accent,below=0.12cm of cic] {$\to f_s/R$};
\node[font=\sffamily\scriptsize,text=accent,below=0.12cm of cfir] {$\to f_s/2R$};
\node[font=\sffamily\scriptsize,text=accent,below=0.12cm of pfir] {$\to f_s/4R$};
\node[font=\sffamily\scriptsize,text=accent,below=0.12cm of rs] {exact rate};
\draw[decorate,decoration={brace,amplitude=5pt},navy!70] ([yshift=4pt]mix.north west) -- ([yshift=4pt]cic.north east)
  node[midway,above=6pt,font=\sffamily\scriptsize,text=navy,align=center]{high rate:\\no multipliers in CIC};
\draw[decorate,decoration={brace,amplitude=5pt},navy!70] ([yshift=4pt]cfir.north west) -- ([yshift=4pt]agc.north east)
  node[midway,above=6pt,font=\sffamily\scriptsize,text=navy]{low rate: sharp filters are cheap here};
\end{tikzpicture}
\caption{A digital down-converter, in the arrangement used by the classic GC4016 and, with half-bands
in place of the CFIR, by the FPGA DDC of the USRP. The expensive arithmetic is pushed to the lowest
sample rates; only the mixer and the CIC integrators run at the ADC rate.}
\label{fig:ch06_ddc}
\end{figure}

The Texas Instruments (originally Graychip) GC4016, a four-channel DDC chip that was a staple of
cellular base stations and software radios around the turn of the century, is the canonical
example: each channel has an NCO and mixer, a five-stage CIC with programmable decimation, a 21-tap
CFIR decimating by two and a 63-tap PFIR decimating by two, followed by a resampler and AGC. Many
later FPGA DDC cores and the DDCs inside RF-sampling converters follow the same pattern, often with
half-band chains in place of the CIC when the decimation is a fixed power of two.

\begin{keyidea}[Push arithmetic to the lowest possible rate]
In a well-designed DDC, only the mixer and the CIC's integrators---a few adders and one complex
multiplication or CORDIC per sample---run at the ADC rate. Every sharp filter runs after most of the
decimation has been done, where it is cheap both because its relative transition is wide and because
it runs rarely. The cost of a DDC is therefore dominated not by the channel filter's sharpness but by
the input rate times the handful of operations that cannot be postponed.
\end{keyidea}

\subsection{Tricks of the trade}

\begin{description}[style=nextline,leftmargin=1.2em,font=\sffamily\bfseries\small\color{navy}]
\item[$f_s/4$ mixing.] If a real IF is placed at exactly a quarter of the sample rate, the mixing
sequence $e^{-j\pi n/2}$ is $1,-j,-1,j,\dots$: the ``mixer'' only routes samples to I or Q and negates
alternate ones, with no multiplier and no NCO. Many IF-sampling receivers choose their IF for this
reason, and finish with a small fine-tuning NCO at a lower rate.
\item[Mixing after decimation.] If the wanted channel lies near the centre of a polyphase decimator's
passband, the frequency shift can be applied to the filter taps instead of the signal (a band-pass
decimator), or after a first coarse decimation, saving high-rate multiplications. GNU Radio's
\texttt{freq\_xlating\_fir\_filter} does exactly this: it rotates the taps to the channel frequency,
filters and decimates, and corrects the residual phase rotation at the output rate.
\item[Real-to-complex in the filter.] A real ADC stream can be converted to a complex one at half the
rate by a complex half-band filter (the $f_s/4$-shifted half-band of Section~\ref{sec:ch06:fir}),
which is also a Hilbert transformer.
\item[Choose the converter clock.] A master clock that is an integer multiple of the signal's natural
rate (30.72 or 61.44\,MHz for LTE and 5G NR, multiples of 1.024\,MHz for DVB-T) removes the need for
fractional resampling altogether.
\end{description}

\subsection{The digital up-converter}

A \emph{digital up-converter} (DUC) is the transpose: a pulse-shaping or channel filter interpolating
by two at the symbol rate, a compensating interpolator by two that pre-distorts for the following
CIC's droop, a CIC interpolator (or half-band chain) up to the DAC rate, and an NCO mixer that places the
channel. Transmitters add two digital stages that receivers lack, crest-factor reduction (clipping and
filtering to lower the peak-to-average ratio) and digital predistortion of the power amplifier
(Chapter~\ref{ch:sdr}), both of which need the sample rate to be several times the signal bandwidth
because they create or cancel out-of-band products. A DAC whose images must be suppressed by an
analog reconstruction filter also benefits from a final inverse-sinc filter that corrects the
DAC's zero-order-hold droop, $\sinc(f/f_s)$, which is 3.9\,dB at $f_s/2$.

\subsection{In software: GNU Radio and SciPy}

In software the same architecture is assembled from library blocks. Table~\ref{tab:ch06:gr} lists the
GNU Radio blocks and SciPy functions corresponding to the techniques of this chapter. The polyphase
blocks are generally the most efficient choice; the FFT filter becomes faster than the direct FIR
for filters of more than a few dozen taps.

\begin{table}[htbp]\centering\small
\caption{Multirate building blocks in GNU Radio (3.10) and SciPy.}\label{tab:ch06:gr}
\begin{tabular}{>{\raggedright\arraybackslash}p{3.0cm}>{\raggedright\arraybackslash}p{6.0cm}>{\raggedright\arraybackslash}p{4.3cm}}
\toprule
Technique & GNU Radio block & SciPy \\
\midrule
FIR design & \texttt{firdes.low\_pass}, \texttt{firdes.root\_raised\_cosine}, \texttt{optfir} & \texttt{firwin}, \texttt{remez}, \texttt{firls} \\
FIR filtering & \texttt{fir\_filter\_ccf}, \texttt{fft\_filter\_ccc} & \texttt{lfilter}, \texttt{oaconvolve} \\
IIR filtering & \texttt{iir\_filter\_ffd}, \texttt{single\_pole\_iir\_filter} & \texttt{sosfilt}, \texttt{sosfiltfilt} \\
Mixing + decimation & \texttt{freq\_xlating\_fir\_filter\_ccc} & --- \\
Polyphase decimation & \texttt{pfb.decimator\_ccf} & \texttt{upfirdn}, \texttt{decimate} \\
Rational resampling & \texttt{rational\_resampler\_ccc} & \texttt{resample\_poly} \\
Arbitrary resampling & \texttt{pfb.arb\_resampler\_ccf} & --- \\
Channelizer & \texttt{pfb.channelizer\_ccf}, \texttt{pfb.synthesizer\_ccf} & --- \\
NCO / mixer & \texttt{sig\_source\_c}, \texttt{rotator\_cc} & --- \\
\bottomrule
\end{tabular}
\end{table}

\begin{pitfall}[Filter state across blocks and calls]
A filter applied to a stream in blocks must carry its state (the last $N-1$ inputs, or the internal
states of each biquad) from one block to the next. Calling \texttt{lfilter} or \texttt{np.convolve}
afresh on each block silently restarts the filter from zero, producing a transient at every block
boundary that appears as a periodic click in audio or a burst of errors in a modem. Pass and return
\texttt{zi}/\texttt{zf} (SciPy) or use a stateful object; and for decimators, also keep track of the
\emph{phase}---which input sample the next output aligns with---or the decimation grid will jump at
each boundary.
\end{pitfall}

%======================================================================
\section{A toolbox of small filters}
\label{sec:ch06:toolbox}
%======================================================================

A few tiny filters appear in almost every radio and deserve to be known by heart.
\begin{description}[style=nextline,leftmargin=1.2em,font=\sffamily\bfseries\small\color{navy}]
\item[The running sum.] The moving average~\eqref{eq:ch06:ma} implemented recursively as
$s[n]=s[n-1]+x[n]-x[n-R]$ costs two additions per sample for any $R$. It computes power estimates,
correlator sums and the energy detectors of burst receivers. In integer arithmetic it is exact; in
floating point the accumulated rounding error must be cleared periodically.
\item[The exponential averager.] $y[n]=y[n-1]+\alpha(x[n]-y[n-1])$, a single pole at $1-\alpha$, is the
cheapest low-pass there is, with a time constant of about $1/\alpha$ samples. With $\alpha=2^{-k}$ it
needs no multiplier. It smooths AGC detectors, RSSI readings and noise-floor estimates.
\item[The DC blocker.] $H(z)=(1-z^{-1})/(1-\alpha z^{-1})$ with $\alpha$ just below 1 is a zero at DC and
a pole just inside it: a high-pass whose $-3$\,dB corner is about $(1-\alpha)f_s/2\pi$. It removes the
residual DC offset of a direct-conversion receiver (Chapter~\ref{ch:sdr}) and is one of the few IIR
filters found in almost every SDR data path.
\item[The Goertzel algorithm.] To compute a single DFT bin $k$ of an $N$-sample block, run the
second-order recursion $s[n]=x[n]+2\cos(2\pi k/N)\,s[n-1]-s[n-2]$ over the block and form
$|X[k]|^2=s^2[N-1]+s^2[N-2]-2\cos(2\pi k/N)\,s[N-1]\,s[N-2]$ at the end. It is a resonator with poles on the unit circle,
used for exactly $N$ samples; it costs one real multiplication per sample per bin, and so beats
an FFT when only a few bins are needed. It is the classic detector for the dual-tone multi-frequency
(DTMF) signalling of telephone keypads (Chapter~\ref{ch:history}).
\end{description}

%======================================================================
\section{Summary}
\label{sec:ch06:summary}
%======================================================================

\begin{itemize}
\item The $z$-transform turns a filter into a pole-zero plot: the frequency response is the ratio of
products of distances from the unit circle to the zeros and to the poles. Poles must lie inside the
circle; a pole at radius $r$ makes a resonance about $(1-r)f_s/\pi$ wide and rings for about
$1/(1-r)$ samples.
\item FIR filter cost is about $A_s/(22\,\Delta f/f_s)$ taps. Window designs (Kaiser) are quick;
Parks--McClellan equiripple designs are optimal in the minimax sense and about 20\% shorter;
least-squares designs minimise stopband energy. Symmetric taps give exact linear phase and halve the
multipliers.
\item Half-band filters have every second tap zero; $M$th-band filters satisfy the Nyquist
criterion; antisymmetric filters make differentiators and Hilbert transformers.
\item IIR filters reach sharp responses at low order but distort phase. Design from analog prototypes
with the bilinear transform and pre-warping; implement as biquads, never as high-order direct forms;
use saturation in recursive paths; beware limit cycles; equalise group delay with all-passes when
needed.
\item In fixed point, accumulate at full precision and round once; scale by $L_1$, $L_\infty$ or $L_2$
norms according to the signal; budget about 6\,dB per bit for coefficients, data and NCO phase.
\item Multirate processing rests on the noble identities and the polyphase decomposition: polyphase
decimators and interpolators never compute discarded outputs or multiply stuffed zeros. Decimate in
stages, with the sharp filter at the lowest rate; half-band chains and CIC filters are the classic
early stages. CIC filters need no multipliers, droop in the passband, and need $N\log_2(RD)$ bits of
register growth, with harmless wrap-around.
\item The polyphase DFT filter bank channelizes $M$ channels for the cost of one filter and one FFT;
the polyphase FFT is a better spectrometer than the windowed FFT.
\item NCOs give phase-continuous, finely resolved tuning; phase truncation costs 6\,dB of SFDR per bit,
and dithering converts spurs to noise. CORDIC rotates and measures vectors with shifts and adds, one
bit per iteration.
\item A DDC chains mixer, CIC, compensating and channel FIRs and a resampler, with only the mixer and
integrators at the ADC rate; the DUC is its transpose.
\end{itemize}

\begin{keyidea}[title=Rules of thumb for radio DSP]
\begin{itemize}
\item Cost of an FIR $\approx A_s/(22\,\Delta f/f_s)$ taps: run sharp filters at low rates.
\item Decimate early and in stages: CIC, then half-bands, then the channel filter.
\item Use polyphase structures so no multiplier ever computes a discarded sample or multiplies a
stuffed zero.
\item Implement IIR filters as biquads; accumulate at full precision and round once.
\item Budget spurs: every bit of NCO phase is 6\,dB, every bit of coefficient 6\,dB of stopband.
\end{itemize}
\end{keyidea}

\begin{labbox}[title=Lab: filters and multirate]
The companion notebook \lab{moderncomms-labs/labs/lab17\_multirate\_dsp.py} (about 90 minutes)
works through this chapter in seven sections. It designs low-pass filters by windowing and by
Parks--McClellan with an interactive tap-count and window selector, and checks Kaiser's length
formula with \texttt{kaiserord}; compares sixth-order Butterworth, Chebyshev and elliptic filters
with a 61-tap FIR, including their group delay; implements a polyphase decimator by 8 and a
polyphase interpolator by 4 and verifies them sample for sample against the direct forms; builds an
integer CIC decimator ($R=16$, $N=4$) with a 31-tap droop compensator and confirms the $R^N$ gain and
register growth; measures the SFDR of an NCO for 6, 10 and 14 phase bits against the 6\,dB-per-bit
rule; and finally assembles a complete DDC---NCO, fixed-point CIC, polyphase compensating FIR and RRC
matched filter---that extracts a 250\,kBd QPSK channel from a 16\,MS/s capture containing strong
neighbours, ending with a clean constellation. Its exercises add a half-band decimate-by-4 chain,
NCO phase dithering and a Farrow resampler. The figure script \lab{book/figscripts/ch06\_figs.py}
generates every data figure in this chapter, including the channelizer, CORDIC and
limit-cycle experiments, and is a good starting point for your own.
\end{labbox}

\problems
\begin{enumerate}[label=\thechapter.\arabic*]
\item A filter has zeros at $z=\pm j$ and poles at $z=0.9e^{\pm j\pi/3}$. Sketch its magnitude response
from the pole-zero geometry, then check it numerically. At what frequency is the gain largest, and
roughly how wide is the peak?
\item Design a notch filter~\eqref{eq:ch06:notch} that removes a 60\,Hz hum from audio sampled at
8\,kHz with a $-3$\,dB width of 2\,Hz. How long does it take to settle? What happens to the notch
depth if the coefficients are rounded to 16 bits in $Q2.14$?
\item Estimate the number of taps for a low-pass FIR with passband to 20\,kHz, stopband from
24\,kHz and 80\,dB attenuation at a sample rate of 192\,kHz using the harris, Bellanger and Kaiser
formulas ($A_p=0.1$\,dB). Repeat at 48\,kHz. Explain the difference.
\item Show that a type II (even-length, symmetric) FIR filter always has a zero at $z=-1$, and that a
type III filter has zeros at $z=\pm1$. Which types can implement a high-pass filter? A Hilbert
transformer? A differentiator?
\item Explain why every other coefficient of a half-band filter is zero. How many multiplications per
input sample does a 31-tap half-band decimator by two need, exploiting both symmetry and the zeros?
Show that shifting it by $f_s/4$ gives a complex filter whose imaginary part is a Hilbert transformer.
\item Use the bilinear transform with pre-warping to design a second-order Butterworth low-pass with a
1\,kHz cutoff at $f_s=8$\,kHz. Find the biquad coefficients, and compute the error in the cutoff if
pre-warping is omitted.
\item Design (in Python) a sixth-order elliptic low-pass filter and implement it in direct form and as
biquads with $Q2.14$ coefficients. Find the smallest word length for which each is stable, and plot the
pole positions of both quantized designs.
\item For the biquad with $a_1=-1.6$, $a_2=0.9$: (a) compute its $L_1$, $L_2$ and $L_\infty$ norms (from
input to output, numerator 1); (b) choose an input scale factor for each scaling rule; (c) for rounding
after each product, estimate the output noise power in LSB$^2$.
\item Prove the noble identity for decimation. Draw the polyphase structure of a decimate-by-3 FIR
filter with 12 taps and count the multiplications per output sample. Do the same for a rational
resampler by $3/2$ with a 24-tap prototype.
\item A CIC decimator has $R=64$, $N=5$ and $D=1$ with a 14-bit input. How many bits must the
integrators have? What is the passband droop at one quarter of the output sample rate? What is the
aliasing attenuation at the inner edge of the first alias band if the wanted signal occupies $\pm0.1$
of the output rate?
\item Design a two-stage decimate-by-40 filter chain for a 200\,kHz channel sampled at 20\,MS/s with
70\,dB of alias protection. Try the splits $20\times2$, $10\times4$, $8\times5$, $5\times8$ and
$4\times10$ and compare their multiplication rates with a single-stage design.
\item An NCO has a 32-bit accumulator clocked at 245.76\,MHz. What is its frequency resolution? What is
the exact output frequency for the tuning word nearest to 10\,MHz, and how fast does its phase drift
relative to an ideal 10\,MHz tone? What phase-word width gives spurs below $-90$\,dBc, and how large
is the quarter-wave table?
\item Derive the CORDIC gain $1/K$ and show that it converges. How many iterations and what datapath
width are needed for 12-bit angle accuracy? What is the range of angles the basic algorithm can
reach?
\item \emph{(Simulation)} Implement CORDIC in vectoring mode in integer arithmetic with 16 stages and
measure its magnitude and phase errors over random inputs. Use it as an FM discriminator (phase
difference between successive samples) and compare with \texttt{np.angle}.
\item \emph{(Simulation)} Build a critically sampled 16-channel polyphase DFT channelizer and verify it
against brute-force mixing, filtering and decimation. Then make it oversampled by two and show that a
carrier placed exactly on a channel boundary can be recovered intact.
\item \emph{(Simulation)} Design a complete DDC from 61.44\,MS/s to 1.92\,MS/s for a 1.4\,MHz LTE carrier
offset by 7\,MHz: NCO, a CIC by 8 with compensator, two half-bands and a 35-tap channel filter. Measure
the EVM of a QPSK test signal with a 20\,dB stronger adjacent carrier, in floating point and with 16-bit
fixed-point data and 18-bit coefficients, and count the multiplications per input sample.
\end{enumerate}

\furtherreading
\begin{itemize}
\item A.~V.~Oppenheim and R.~W.~Schafer, \emph{Discrete-Time Signal Processing}, 3rd ed., Pearson,
2010. The standard graduate text: the $z$-transform, filter structures, quantization effects and
multirate basics, all derived carefully.
\item R.~G.~Lyons, \emph{Understanding Digital Signal Processing}, 3rd ed., Pearson, 2010. Exceptionally
clear, and full of implementation tricks---CIC filters, DC blockers, Goertzel, sharpening---that the
formal texts omit.
\item f.~j.~harris, \emph{Multirate Signal Processing for Communication Systems}, 2nd ed., River
Publishers, 2021. The practitioner's bible for polyphase filters, channelizers, arbitrary resamplers
and half-band designs, written by the engineer who built many of them.
\item R.~E.~Crochiere and L.~R.~Rabiner, \emph{Multirate Digital Signal Processing}, Prentice Hall,
1983; and their tutorial ``Interpolation and decimation of digital signals---a tutorial review,''
\emph{Proc. IEEE}, vol.~69, no.~3, 1981. The foundations of multistage design.
\item P.~P.~Vaidyanathan, \emph{Multirate Systems and Filter Banks}, Prentice Hall, 1993. The definitive
treatment of polyphase theory, perfect-reconstruction filter banks and their design.
\item E.~B.~Hogenauer, ``An economical class of digital filters for decimation and interpolation,''
\emph{IEEE Trans. Acoustics, Speech, and Signal Processing}, vol.~29, no.~2, pp.~155--162, 1981. The CIC
paper, including register growth and pruning; short and very readable.
\item T.~W.~Parks and J.~H.~McClellan, ``Chebyshev approximation for nonrecursive digital filters with
linear phase,'' \emph{IEEE Trans. Circuit Theory}, vol.~19, no.~2, 1972; and J.~H.~McClellan,
T.~W.~Parks and L.~R.~Rabiner, ``A computer program for designing optimum FIR linear phase digital
filters,'' \emph{IEEE Trans. Audio and Electroacoustics}, vol.~21, no.~6, 1973.
\item M.~G.~Bellanger, G.~Bonnerot and M.~Coudreuse, ``Digital filtering by polyphase network:
application to sample-rate alteration and filter banks,'' \emph{IEEE Trans. Acoustics, Speech, and
Signal Processing}, vol.~24, no.~2, 1976. The origin of the polyphase decomposition.
\item J.~E.~Volder, ``The CORDIC trigonometric computing technique,'' \emph{IRE Trans. Electronic
Computers}, vol.~EC-8, no.~3, 1959; and R.~Andraka, ``A survey of CORDIC algorithms for FPGA based
computers,'' \emph{Proc. ACM/SIGDA FPGA '98}, 1998, the standard hardware-oriented introduction.
\item L.~B.~Jackson, \emph{Digital Filters and Signal Processing}, 3rd ed., Kluwer, 1996. Especially good
on fixed-point effects: scaling, roundoff noise, limit cycles and section ordering in IIR cascades.
\item Analog Devices, \emph{AD9361 Reference Manual} (UG-570). Describes the transceiver's half-band
and programmable FIR chain and its design tools in detail; read alongside the UHD documentation for
the B200's FPGA DDC.
\end{itemize}

